\documentclass[11pt]{article}
\usepackage[margin=1in]{geometry}
\usepackage[T1]{fontenc}
\usepackage[utf8]{inputenc}
\usepackage[french]{babel}
\usepackage{amsmath,amssymb,amsthm}
\usepackage[colorlinks=true,linkcolor=blue,urlcolor=blue]{hyperref}
\newtheorem{problem}{Problème}
\newenvironment{refs}{\par\small\noindent\emph{Références :}\begin{itemize}\setlength\itemsep{0pt}}{\end{itemize}\normalsize}
\title{Hors de la Caverne \\ \Large L'art de la démonstration}
\author{}
\date{}
\begin{document}
\maketitle
\noindent\emph{Des problèmes, pas de réponses.} Chaque problème ci-dessous est posé
sans solution. Les références indiquent où trouver les connaissances nécessaires.
\medskip


\section*{Collège}

\begin{problem}[{\og{} J'ai essayé avec 3, 5 et 7, donc c'est vrai \fg{} --- Collège}]
\textbf{PRF-64.} \textbf{Problème.} Un camarade affirme : \og{} J'ai essayé avec 3, 5 et 7, ça marche à chaque
fois, donc c'est vrai. \fg{} Ce problème demande ce qui manque exactement à cette phrase.
\begin{enumerate}
\item[(a)] Calculez $ n^2 + n + 11 $ pour n = 0, 1, 2, \dots{}, 9, et constatez que vous obtenez chaque
fois un nombre premier. Décidez ensuite si l'affirmation \og{} $ n^2 + n + 11 $ est un nombre
premier pour tout entier naturel n \fg{} est vraie, et démontrez votre réponse.
\item[(b)] Placez 2 points sur un cercle et tracez la corde qui les joint : le disque est coupé en
2 morceaux. Recommencez avec 3, puis 4, puis 5, puis 6 points, en traçant toutes les cordes
et en plaçant les points de sorte que jamais trois cordes ne se croisent au même endroit ;
comptez les morceaux à chaque fois. Écrivez la suite de nombres obtenue, dites quelle règle
elle vous suggère, puis dites si vos comptages la démontrent.
\item[(c)] Un seul contre-exemple suffit-il à réfuter une affirmation qui commence par \og{} pour
tout \fg{} ? Suffit-il à réfuter une affirmation qui commence par \og{} il existe \fg{} ? Justifiez les
deux réponses.
\item[(d)] Rédigez en trois ou quatre phrases ce qu'il faudrait ajouter à \og{} j'ai essayé avec 3, 5
et 7 \fg{} pour en faire une démonstration --- et pourquoi mille essais de plus n'y suffiraient
toujours pas.
\end{enumerate}

\end{problem}

\noindent Les deux pièges de ce problème sont célèbres et ils ont piégé de vrais
mathématiciens. Euler remarqua en 1772 que $ n^2 + n + 41 $ donne un nombre premier pour
quarante valeurs consécutives de n avant de céder ; 11 appartient à la même petite famille
de nombres, dits nombres chanceux d'Euler, et c'est pour cela qu'il tient si longtemps.
Le compte des morceaux du disque est le piège inverse : la suite commence si franchement
qu'on croit la connaître, et le sixième terme n'est pas celui qu'on annonce --- Leo Moser en
a fait un classique de la mise en garde. La leçon est la même dans les deux cas et elle est
au fondement de tout ce site : vérifier n'est pas démontrer, l'un donne de la confiance,
l'autre donne la certitude, et rien ne transforme la première en la seconde.

\begin{refs}
\item Contexte : Contre-exemple --- Wikipédia (\url{https://fr.wikipedia.org/wiki/Contre-exemple})
\item Contexte : Nombre chanceux d'Euler --- Wikipédia (\url{https://fr.wikipedia.org/wiki/Nombre_chanceux_d%27Euler})
\item Contexte : Partage d'un disque par des cordes --- Wikipédia (en anglais) (\url{https://en.wikipedia.org/wiki/Dividing_a_circle_into_areas})
\end{refs}

\bigskip

\begin{problem}[{Deux impairs font un pair --- démontrez-le --- Collège}]
\textbf{PRF-65.} \textbf{Problème.} Un entier est pair s'il s'écrit $ 2k $ pour un certain entier k, et
impair s'il s'écrit $ 2k+1 $ pour un certain entier k. Tout ce qui suit ne doit se servir
que de ces deux écritures.
\begin{enumerate}
\item[(a)] Vérifiez sur cinq couples de votre choix que la somme de deux entiers impairs est
paire, puis dites en une phrase pourquoi ces cinq vérifications ne démontrent rien.
\item[(b)] Démontrez que la somme de deux entiers impairs est paire. Votre argument doit valoir
pour tous les couples d'entiers impairs d'un seul coup ; montrez du doigt l'endroit précis
où il cesse de parler d'exemples.
\item[(c)] Justifiez le point sur lequel repose (b) et qu'on oublie toujours de justifier : tout
entier impair s'écrit bien $ 2k+1 $ pour un certain entier k, et tout nombre de cette
forme est bien impair.
\item[(d)] Démontrez de la même façon que le produit de deux entiers impairs est impair, que la
somme de trois entiers impairs est impaire, et que la somme de trois entiers consécutifs est
divisible par 3. Décidez enfin si la somme de quatre entiers consécutifs est divisible
par 4.
\end{enumerate}

\end{problem}

\noindent C'est presque toujours la première vraie démonstration qu'un élève écrit,
et elle contient déjà tout le mécanisme : une lettre remplace un nombre, et la ligne de
calcul qui suit parle de tous les cas à la fois. Les Grecs travaillaient ces énoncés bien
avant qu'on sache écrire une lettre pour un nombre --- le livre IX des \textit{Éléments}
d'Euclide démontre en propositions séparées que la somme de nombres pairs est paire, que la
somme d'un nombre pair de nombres impairs est paire, et ainsi de suite, entièrement en
langue et en longueurs. Comparer leur rédaction à la vôtre est instructif : le calcul
littéral n'a pas rendu la démonstration possible, il l'a rendue courte. PRF-74 vous fera
revenir sur le texte que vous aurez écrit ici.

\begin{refs}
\item Contexte : Parité (arithmétique) --- Wikipédia (\url{https://fr.wikipedia.org/wiki/Parit%C3%A9_(arithm%C3%A9tique)})
\item Contexte : Éléments d'Euclide --- Wikipédia (\url{https://fr.wikipedia.org/wiki/%C3%89l%C3%A9ments_d%27Euclide})
\item Contexte : Calcul algébrique --- Wikipédia (\url{https://fr.wikipedia.org/wiki/Calcul_alg%C3%A9brique})
\end{refs}

\bigskip

\begin{problem}[{Un nombre et son double --- Collège, short}]
\textbf{PRF-66.} \textbf{Problème.} Démontrez ou réfutez : un entier et son double n'ont
jamais la même parité. Si l'affirmation est fausse, énoncez à sa place la règle exacte qui
dit quand un entier n et son double $ 2n $ ont la même parité, et démontrez cette
règle.
\end{problem}

\noindent \og{} Démontrez ou réfutez \fg{} est un genre à part : la moitié du travail
consiste à décider de quel côté on se range, et la première chose à faire est d'essayer de
casser l'énoncé plutôt que de le croire. Beaucoup d'affirmations fausses en mathématiques
sont presque vraies --- vraies sur la moitié des cas, ou sur tous sauf un --- et c'est
précisément ce qui les rend dangereuses. PRF-38 pose la même question au lycée, sur les
nombres irrationnels.

\begin{refs}
\item Contexte : Parité (arithmétique) --- Wikipédia (\url{https://fr.wikipedia.org/wiki/Parit%C3%A9_(arithm%C3%A9tique)})
\item Contexte : Contre-exemple --- Wikipédia (\url{https://fr.wikipedia.org/wiki/Contre-exemple})
\end{refs}

\bigskip

\begin{problem}[{Où se cache l'erreur : une \og{} démonstration \fg{} que 1 = 2 --- Collège, short}]
\textbf{PRF-67.} \textbf{Problème.} On part de deux nombres égaux, $ a = b $, et on
enchaîne :
\[ a^2 = ab, \qquad a^2 - b^2 = ab - b^2, \qquad (a-b)(a+b) = b(a-b), \qquad a+b = b, \qquad 2b = b, \qquad 2 = 1. \]
Désignez la ligne exactement fausse --- pas \og{} quelque part par là \fg{}, la ligne --- et écrivez en
une phrase la règle qu'elle enfreint.
\end{problem}

\noindent Toutes les lignes sont correctement calculées, et pourtant la conclusion
est fausse : c'est le meilleur argument qui soit contre l'idée qu'une démonstration serait
un calcul juste. Une démonstration est un enchaînement dont chaque pas est \textit{autorisé},
et une seule autorisation manquante suffit à tout ruiner. Ce faux raisonnement circule
depuis des siècles et il a un cousin géométrique tout aussi instructif, le paradoxe du carré
manquant, où un découpage semble faire apparaître une case de plus : là, ce qui triche n'est
pas un calcul mais un dessin. Retenez la morale des deux --- une figure et une ligne de calcul
sont des indices, jamais des preuves.

\begin{refs}
\item Contexte : Division par zéro --- Wikipédia (\url{https://fr.wikipedia.org/wiki/Division_par_z%C3%A9ro})
\item Contexte : Raisonnement fallacieux --- Wikipédia (\url{https://fr.wikipedia.org/wiki/Raisonnement_fallacieux})
\item Contexte : Paradoxe du carré manquant --- Wikipédia (\url{https://fr.wikipedia.org/wiki/Paradoxe_du_carr%C3%A9_manquant})
\end{refs}

\bigskip

\begin{problem}[{Le principe des tiroirs, raconté puis employé --- Collège}]
\textbf{PRF-68.} \textbf{Problème.} Le principe des tiroirs dit ceci : si l'on range plus d'objets qu'on n'a
de tiroirs, alors un tiroir au moins contient au moins deux objets. Il tient en une ligne et
il a l'air de ne rien dire --- pourtant tout ce qui suit en découle.
\begin{enumerate}
\item[(a)] Démontrez le principe lui-même, en deux phrases, par l'absurde : supposez que chaque
tiroir contienne au plus un objet, et comptez.
\item[(b)] Démontrez que parmi 13 personnes quelconques, deux au moins sont nées le même mois.
\item[(c)] On choisit six nombres différents parmi 1, 2, \dots{}, 10. Démontrez que deux d'entre eux ont
pour somme 11.
\item[(d)] Une tête humaine porte au plus 150 000 cheveux. Démontrez que dans une ville de
200 000 habitants, deux personnes au moins ont exactement le même nombre de cheveux.
\item[(e)] Voici le véritable exercice. Pour chacune des questions (b), (c) et (d), écrivez
explicitement ce qui jouait le rôle des objets, ce qui jouait celui des tiroirs, et combien
il y en avait de chaque --- puis dites, en une phrase par question, lequel des deux était le
plus difficile à voir.
\end{enumerate}

\end{problem}

\noindent Le principe porte le nom de Dirichlet, qui s'en est servi comme d'un outil
en 1834 sous le nom de \textit{Schubfachprinzip}, \og{} principe des tiroirs \fg{}, mais l'exemple des
cheveux est bien plus ancien : un recueil français de 1622, la \textit{Récréation
mathématique} publiée sous le nom de Jean Leurechon, remarque déjà qu'il est nécessaire
que deux hommes aient le même nombre de cheveux. Ce qui frappe, c'est le contraste entre la
banalité de l'énoncé et la force de ses conséquences : il produit des affirmations qu'aucune
observation ne pourrait établir, sur des gens qu'on n'a jamais vus. Toute la difficulté, à
chaque emploi, tient dans la question \og{} quels sont les tiroirs ? \fg{} --- car ils ne sont jamais
donnés. PRF-07 et PRF-55 reprennent la question au lycée, avec des tiroirs qu'il faut
fabriquer soi-même.

\begin{refs}
\item Contexte : Principe des tiroirs --- Wikipédia (\url{https://fr.wikipedia.org/wiki/Principe_des_tiroirs})
\item Contexte : Peter Gustav Lejeune Dirichlet --- Wikipédia (\url{https://fr.wikipedia.org/wiki/Peter_Gustav_Lejeune_Dirichlet})
\item Lecture : A. Engel, \textit{Problem-Solving Strategies}, ch. 4 (le principe des tiroirs)
\end{refs}

\bigskip

\begin{problem}[{Deux personnes ont le même nombre d'amis --- Collège, short}]
\textbf{PRF-69.} \textbf{Problème.} Dans un groupe d'au moins deux personnes, où l'amitié
est toujours réciproque, démontrez qu'il existe deux personnes ayant exactement le même
nombre d'amis à l'intérieur du groupe. Rédigez-le comme un raisonnement par l'absurde :
supposez que ces nombres soient tous différents, et poussez cette supposition jusqu'à une
impossibilité.
\end{problem}

\noindent Le raisonnement par l'absurde est la manœuvre la plus étrange de la
logique --- on affirme le contraire de ce qu'on veut établir, on travaille consciencieusement
sur une hypothèse dont on espère qu'elle est fausse, et l'effondrement final est la
démonstration. Ce problème-ci en est une version où l'on ne calcule rien du tout : il n'y a
ni formule, ni figure, et l'énoncé ne dit même pas combien de personnes il y a. Traduit en
langage de graphes --- un point par personne, un trait par amitié --- c'est l'un des premiers
résultats qu'on rencontre dans la théorie née de l'article d'Euler sur les ponts de
Königsberg (1736).

\begin{refs}
\item Contexte : Raisonnement par l'absurde --- Wikipédia (\url{https://fr.wikipedia.org/wiki/Raisonnement_par_l%27absurde})
\item Contexte : Théorie des graphes --- Wikipédia (\url{https://fr.wikipedia.org/wiki/Th%C3%A9orie_des_graphes})
\item Contexte : Principe des tiroirs --- Wikipédia (\url{https://fr.wikipedia.org/wiki/Principe_des_tiroirs})
\end{refs}

\bigskip

\begin{problem}[{Ce qui ne change jamais --- Collège}]
\textbf{PRF-70.} \textbf{Problème.} Un invariant est une quantité qui ne bouge pas alors que la situation,
elle, bouge. Chacune des questions ci-dessous demande de démontrer qu'une chose est
\textit{impossible} --- et aucun essai raté, si nombreux soient-ils, ne peut établir cela.
\begin{enumerate}
\item[(a)] Sept verres sont posés à l'envers sur une table. Un coup consiste à retourner exactement
deux verres, lesquels que vous voulez. Démontrez qu'aucune suite de coups ne mettra les sept
verres à l'endroit.
\item[(b)] On écrit au tableau les nombres 1, 2, \dots{}, 10. Un coup consiste à effacer deux nombres et
à écrire à la place leur différence positive (ou 0 s'ils sont égaux). Après neuf coups il ne
reste qu'un nombre : démontrez que sa parité est fixée d'avance, quels qu'aient été vos
choix.
\item[(c)] Un jeton est posé sur une case d'un quadrillage colorié comme un damier. À chaque coup
il se déplace en diagonale, d'une case exactement. Démontrez qu'il ne pourra jamais atteindre
une case de l'autre couleur, et dites ce que cela raconte du fou aux échecs.
\item[(d)] Pour chacune des trois questions, nommez l'invariant en une phrase. Expliquez ensuite
pourquoi exhiber un invariant démontre une impossibilité, alors que cent tentatives ratées
ne démontrent rien du tout.
\end{enumerate}

\end{problem}

\noindent L'argument d'invariance est peut-être la plus belle idée élémentaire des
mathématiques : au lieu de suivre un processus qui part dans toutes les directions, on
trouve la seule chose qu'il ne peut pas modifier, et l'on regarde si le but demandé la
respecte. S'il ne la respecte pas, le but est hors d'atteinte --- définitivement, pour toutes
les suites de coups possibles, y compris celles que personne n'essaiera jamais. C'est la
première technique du livre d'Arthur Engel pour les olympiades, et ce n'est pas un hasard :
c'est aussi la manière dont on démontre qu'un programme informatique fait bien ce qu'il
prétend (DCS-01 en donne la version échiquier). Le mot \og{} impossible \fg{} est l'un des rares que
les mathématiques prononcent sans trembler ; l'invariant est ce qui leur en donne le
droit.

\begin{refs}
\item Contexte : Invariant --- Wikipédia (\url{https://fr.wikipedia.org/wiki/Invariant})
\item Lecture : A. Engel, \textit{Problem-Solving Strategies}, ch. 1 (le principe d'invariance)
\item Voir aussi : DCS-01, l'échiquier mutilé (\url{applied-discrete-cs.html#DCS-01})
\end{refs}

\bigskip

\begin{problem}[{Une case en moins, et le coloriage qui décide --- Collège}]
\textbf{PRF-71.} \textbf{Problème.} On colorie un carré $ 5 \times 5 $ comme un damier, les quatre coins
recevant la même couleur. On retire une case, puis on essaie de recouvrir les 24 cases
restantes par 12 dominos, chaque domino couvrant exactement deux cases voisines par un
côté.
\begin{enumerate}
\item[(a)] Comptez les cases de chaque couleur. Démontrez ensuite qu'aucun pavage n'existe dès que
la case retirée est de la couleur la moins nombreuse.
\item[(b)] Cherchez maintenant, parmi les cases de l'autre couleur, celles pour lesquelles un
pavage existe réellement. Chaque réponse demande sa justification : une construction dans un
sens, une démonstration d'impossibilité dans l'autre.
\item[(c)] Le coloriage ne figure nulle part dans l'énoncé du problème : c'est vous qui l'ajoutez.
Expliquez en quelques phrases ce qu'il rend visible et que le quadrillage nu ne montrait
pas.
\item[(d)] Reprenez (a) en supposant qu'un domino ait aussi le droit de couvrir deux cases voisines
\textit{en diagonale}. Dites ce que devient l'argument, et à quelle ligne exactement il
cesse de fonctionner.
\end{enumerate}

\end{problem}

\noindent Colorier une figure pour démontrer qu'un découpage est impossible est un
tour de main que Solomon Golomb a installé au cœur d'un domaine entier lorsqu'il a inventé
le mot \og{} polyomino \fg{} en 1954 et publié ses premiers résultats de pavage. La force du procédé
tient à ce qu'il ajoute au problème une information qui n'y était pas : le damier n'existe
que dans votre tête, et pourtant c'est lui qui tranche. Remarquez aussi ce que la question
(b) demande de plus que la (a) --- une impossibilité se démontre par un argument, mais une
possibilité se démontre en montrant la chose, et les deux moitiés d'un \og{} pour quelles cases
est-ce faisable ? \fg{} n'ont jamais la même allure. Le résultat frère sur l'échiquier
$ 8 \times 8 $ privé de deux cases de couleurs différentes est dû à Ralph Gomory, et son
argument --- d'une élégance rare --- a été popularisé par Ross Honsberger en 1973.

\begin{refs}
\item Contexte : Polyomino --- Wikipédia (\url{https://fr.wikipedia.org/wiki/Polyomino})
\item Contexte : Problème de l'échiquier mutilé --- Wikipédia (\url{https://fr.wikipedia.org/wiki/Probl%C3%A8me_de_l%27%C3%A9chiquier_mutil%C3%A9})
\item Lecture : S. W. Golomb, \textit{Polyominoes} (2ᵉ éd., 1994), ch. 1
\end{refs}

\bigskip

\begin{problem}[{La réciproque n'est pas l'énoncé --- Collège}]
\textbf{PRF-72.} \textbf{Problème.} La réciproque de \og{} si P alors Q \fg{} est \og{} si Q alors P \fg{}. Ce sont deux
affirmations différentes : démontrer l'une ne démontre pas l'autre, et il arrive
constamment que l'une soit vraie et l'autre fausse.
\begin{enumerate}
\item[(a)] Pour chacun des énoncés suivants, écrivez la réciproque, puis décidez si cette
réciproque est vraie --- par une démonstration si oui, par un contre-exemple si non. (i) Si un
entier se termine par 0, alors il est divisible par 5. (ii) Si un triangle est équilatéral,
alors il est isocèle. (iii) Si un entier est divisible par 6, alors il est divisible par 2
et par 3. (iv) S'il pleut, le sol est mouillé.
\item[(b)] Ceux des énoncés dont la réciproque est vraie méritent la forme \og{} si et seulement si \fg{}.
Réécrivez-les ainsi, et rédigez pour l'un d'eux les deux sens comme deux démonstrations
clairement séparées et étiquetées.
\item[(c)] Un critère de divisibilité dit : si la somme des chiffres d'un entier est divisible par
3, alors l'entier l'est aussi. Écrivez sa réciproque, et dites laquelle des deux vous
employez réellement quand vous testez 1 234 567.
\item[(d)] Expliquez en trois phrases pourquoi \og{} la réciproque est évidente \fg{} est l'une des phrases
les plus dangereuses qu'on puisse écrire dans une copie de mathématiques.
\end{enumerate}

\end{problem}

\noindent Confondre un énoncé et sa réciproque est l'erreur de raisonnement la plus
répandue, et pas seulement en mathématiques : \og{} tous les menteurs se contredisent \fg{} ne dit
rien de celui qui vient de se contredire, et un tribunal qui l'oublie condamne un innocent.
Aristote la traitait déjà comme une faute nommée, et la logique médiévale l'appelait
l'affirmation du conséquent. En classe, l'occasion la plus fréquente de rencontrer la
distinction est géométrique : le théorème de Thalès sert à calculer des longueurs, sa
réciproque sert à démontrer que des droites sont parallèles, et ce sont deux outils
différents (PRF-73). PRF-46 reprend la question au lycée, sur des équivalences plus
coriaces.

\begin{refs}
\item Contexte : Implication réciproque --- Wikipédia (\url{https://fr.wikipedia.org/wiki/Implication_r%C3%A9ciproque})
\item Contexte : Implication (logique) --- Wikipédia (\url{https://fr.wikipedia.org/wiki/Implication_(logique)})
\item Contexte : Critères de divisibilité --- Wikipédia (\url{https://fr.wikipedia.org/wiki/Crit%C3%A8res_de_divisibilit%C3%A9})
\end{refs}

\bigskip

\begin{problem}[{Thalès, et la réciproque de Thalès --- Collège}]
\textbf{PRF-73.} \textbf{Problème.} Deux droites se coupent en A. Sur la première on place B, puis M, du même
côté de A ; sur la seconde on place C, puis N, du même côté de A. Le théorème de Thalès
affirme que si les droites $ (BC) $ et $ (MN) $ sont parallèles, alors
\[ \frac{AB}{AM} \;=\; \frac{AC}{AN} \;=\; \frac{BC}{MN}. \]
\begin{enumerate}
\item[(a)] On donne $ AB = 3 $, $ AM = 5 $, $ AC = 6 $, et $ (BC) $ parallèle à
$ (MN) $. Calculez AN en écrivant, à côté de chaque égalité, l'hypothèse qui l'autorise.
Une seule de vos lignes se sert du parallélisme : entourez-la.
\item[(b)] La réciproque du théorème sert à démontrer qu'une chose est parallèle. Énoncez-la
précisément, hypothèses comprises, puis servez-vous-en pour démontrer que deux droites d'une
figure de votre construction sont parallèles.
\item[(c)] La réciproque exige que les points soient placés \og{} dans le même ordre \fg{} sur les deux
droites. Construisez une figure où les rapports de longueurs sont bien égaux mais où cette
condition n'est pas remplie, et décidez, avec justification, si les droites y sont
parallèles.
\item[(d)] Dites en deux phrases ce que (c) enseigne : ce qui arrive à une démonstration quand on
applique un théorème sans vérifier l'une de ses hypothèses.
\end{enumerate}

\end{problem}

\noindent La tradition fait de Thalès de Milet, vers 600 av. J.-C., le premier
à avoir \textit{démontré} plutôt que constaté, et la plus jolie des légendes le montre mesurant
la hauteur de la grande pyramide par la longueur de son ombre --- c'est-à-dire employant très
exactement le théorème qui porte aujourd'hui son nom. Une curiosité mérite d'être connue :
l'appellation \og{} théorème de Thalès \fg{} pour cet énoncé de proportions est une convention
française, quand l'anglais l'appelle le théorème des intercepts et réserve le nom de Thalès
à un tout autre résultat, celui de l'angle inscrit dans un demi-cercle. Ce théorème est
souvent le premier dont un élève français utilise la réciproque comme un outil de
démonstration à part entière, et il vaut la peine de bien voir que les deux sens ne servent
pas au même travail : l'un mesure, l'autre prouve.

\begin{refs}
\item Contexte : Théorème de Thalès --- Wikipédia (\url{https://fr.wikipedia.org/wiki/Th%C3%A9or%C3%A8me_de_Thal%C3%A8s})
\item Contexte : Thalès --- Wikipédia (\url{https://fr.wikipedia.org/wiki/Thal%C3%A8s})
\item Contexte : Éléments d'Euclide --- Wikipédia (\url{https://fr.wikipedia.org/wiki/%C3%89l%C3%A9ments_d%27Euclide})
\end{refs}

\bigskip

\begin{problem}[{Rédiger : ce qu'un lecteur a le droit d'exiger --- Collège}]
\textbf{PRF-74.} \textbf{Problème.} Trouver un argument et l'écrire sont deux travaux distincts, et le second
s'apprend. Reprenez la démonstration que vous avez trouvée en PRF-65 (b) --- la somme de deux
entiers impairs est paire --- et travaillez-la comme on travaille un texte.
\begin{enumerate}
\item[(a)] Rédigez-la pour un camarade absent ce jour-là : chaque fait employé doit être nommé, et
aucune étape ne doit lui demander de deviner. Comptez ensuite vos phrases.
\item[(b)] Soulignez les mots qui portent la logique --- \og{} soit \fg{}, \og{} supposons \fg{}, \og{} donc \fg{}, \og{} or \fg{},
\og{} il existe \fg{}, \og{} pour tout \fg{} --- et vérifiez qu'aucune phrase n'en est privée par accident.
Une phrase sans mot de liaison est souvent une phrase qui ne démontre rien.
\item[(c)] Donnez votre texte à quelqu'un qui n'a pas cherché le problème, en lui demandant de
s'arrêter à la première ligne qu'il ne croit pas. Réécrivez à partir de cette ligne, puis
recommencez l'opération avec un autre lecteur.
\item[(d)] Faites le même travail sur votre démonstration de PRF-68 (c), qui est d'une autre
nature : une phrase y remplace un calcul. Dites enfin, en deux phrases, ce que change le
passage de \og{} je vois pourquoi c'est vrai \fg{} à \og{} j'ai écrit pourquoi c'est vrai \fg{}.
\end{enumerate}

\end{problem}

\noindent La quatrième et dernière étape de la méthode de George Pólya, dans
\textit{How to Solve It} (1945), traduit en français sous le titre \textit{Comment poser et
résoudre un problème}, s'appelle \og{} revenir en arrière \fg{} : le problème n'est pas fini quand
la réponse est trouvée, il est fini quand l'argument tient debout tout seul, sans son auteur
à côté pour l'expliquer. C'est exactement le service que rend, dans la vie d'un
mathématicien, le rapporteur d'un article : quelqu'un qui n'est pas dans votre tête et qui
s'arrête à la première ligne qu'il ne croit pas. La question (c) vous fait fabriquer ce
lecteur, et c'est de loin l'exercice le plus utile de cette page --- parce qu'une démonstration
qu'on est seul à comprendre n'a encore convaincu personne.

\begin{refs}
\item Contexte : Démonstration (logique et mathématiques) --- Wikipédia (\url{https://fr.wikipedia.org/wiki/D%C3%A9monstration_(logique_et_math%C3%A9matiques)})
\item Contexte : George Pólya --- Wikipédia (\url{https://fr.wikipedia.org/wiki/George_P%C3%B3lya})
\item Lecture : G. Pólya, \textit{Comment poser et résoudre un problème} (trad. fr. de \textit{How to Solve It})
\end{refs}

\bigskip

\begin{problem}[{Toujours vrai, ou seulement souvent vrai --- Collège, short}]
\textbf{PRF-75.} \textbf{Problème.} Décidez, avec démonstration, si l'affirmation \og{} la
somme de deux nombres premiers est un nombre pair \fg{} est toujours vraie ; si elle ne l'est
pas, énoncez et démontrez la règle exacte qui dit quand elle l'est. Dites ensuite en une
phrase ce que change, pour un mathématicien, le passage de \og{} je n'ai trouvé aucune
exception \fg{} à \og{} il ne peut pas y en avoir \fg{}.
\end{problem}

\noindent La distinction entre \og{} toujours vrai \fg{} et \og{} vrai dans tous les cas que
j'ai regardés \fg{} est la frontière même des mathématiques, et il existe un exemple magnifique
juste à côté de celui-ci : la conjecture de Goldbach affirme que tout entier pair supérieur
à 2 est la somme de deux nombres premiers. Elle a été vérifiée par ordinateur pour tous les
entiers pairs jusqu'à plus de quatre milliards de milliards, sans la moindre exception, et
personne ne l'a démontrée depuis qu'Euler et Goldbach en ont discuté par lettres en 1742.
Voilà à quoi ressemble une affirmation \og{} souvent vraie \fg{} : indéfiniment vérifiée, toujours
non démontrée, et donc toujours pas un théorème.

\begin{refs}
\item Contexte : Nombre premier --- Wikipédia (\url{https://fr.wikipedia.org/wiki/Nombre_premier})
\item Contexte : Conjecture de Goldbach --- Wikipédia (\url{https://fr.wikipedia.org/wiki/Conjecture_de_Goldbach})
\item Voir aussi : Problèmes ouverts (\url{open-problems.html})
\end{refs}

\bigskip


\section*{Lycée}

\begin{problem}[{Démonstration directe ou contraposée ? --- Lycée}]
\textbf{PRF-01.} \textbf{Problème.} Une implication \og{} si P alors Q \fg{} peut se démontrer directement
(on suppose P, on en déduit Q) ou par contraposée (on suppose Q fausse, on en déduit que
P est fausse). Pour chacun des énoncés ci-dessous, décidez laquelle des deux formes donne
la démonstration la plus nette, et rédigez-la.
\begin{enumerate}
\item[(a)] Si n est un entier impair, alors $ n^2 $ est impair.
\item[(b)] Si $ n^2 $ est pair pour un entier n, alors n est pair.
\item[(c)] Si x et y sont des réels strictement positifs avec $ xy >25 $, alors $ x >5 $ ou
$ y >5 $.
\item[(d)] Si la somme de deux entiers vaut au moins 100, alors l'un au moins des deux vaut au
moins 50.
\item[(e)] Formulez ensuite, en une ou deux phrases, une règle empirique pour reconnaître quand
la contraposée sera la voie la plus facile.
\end{enumerate}

\end{problem}

\noindent Qu'une implication et sa contraposée soient vraies ou fausses ensemble
était déjà clair pour Aristote, et ce fut un lieu commun de la logique médiévale
(\textit{modus tollens}). Choisir la bonne forme est la première décision véritable que prend
celui qui rédige une démonstration : un énoncé dont la conclusion contient un \og{} ou \fg{}, ou
dont l'hypothèse est difficile à saisir, devient souvent docile dès qu'on le retourne.
Toute introduction à la démonstration commence ici.

\begin{refs}
\item Contexte : Proposition contraposée --- Wikipédia (\url{https://fr.wikipedia.org/wiki/Proposition_contrapos%C3%A9e})
\item Lecture : D. J. Velleman, \textit{How to Prove It}, ch. 3
\item Lecture : R. Hammack, \textit{Book of Proof}, ch. 4--5 (gratuit en ligne (\url{https://richardhammack.github.io/BookOfProof/}))
\end{refs}

\bigskip

\begin{problem}[{La racine carrée de 2 est irrationnelle --- Lycée}]
\textbf{PRF-02.} \textbf{Problème.} Démontrez que $ \sqrt{2} $ est irrationnel : il
n'existe pas d'entiers strictement positifs p et q tels que $ \sqrt{2} = p/q $. (Forme
suggérée : le raisonnement par l'absurde --- supposez qu'une telle fraction existe et partez
en chasse d'une impossibilité.)
\end{problem}

\noindent Le plus vieux coup de tonnerre des mathématiques. Les pythagoriciens, pour
qui \og{} tout est nombre \fg{} voulait dire les entiers et leurs rapports, ont découvert que la
diagonale d'un carré est incommensurable avec son côté --- la légende veut que la découverte
ait été imputée à Hippase de Métaponte, noyé en mer pour l'avoir révélée. Aristote cite
l'argument pair-impair comme exemple type du raisonnement par l'impossible, et une version
en clôt le livre X des \textit{Éléments} d'Euclide. C'est le modèle du raisonnement par
l'absurde, et le problème suivant demande de quoi ce modèle est réellement fait.

\begin{refs}
\item Contexte : Racine carrée de deux --- Wikipédia (\url{https://fr.wikipedia.org/wiki/Racine_carr%C3%A9e_de_deux})
\item Lecture : G. H. Hardy \& E. M. Wright, \textit{An Introduction to the Theory of Numbers}, ch. 4
\item Lecture : R. Hammack, \textit{Book of Proof}, ch. 6
\end{refs}

\bigskip

\begin{problem}[{Deux autres irrationnels, et la forme du canevas --- Lycée}]
\textbf{PRF-03.} \textbf{Problème.} PRF-02 a démontré que $ \sqrt{2} $ est irrationnel ; l'argument
est un canevas réutilisable, et ce problème vous fait l'en extraire.
\begin{enumerate}
\item[(a)] Démontrez que $ \sqrt{3} $ est irrationnel.
\item[(b)] Démontrez que $ \log_2 3 $ est irrationnel.
\item[(c)] Prenez maintenant du recul : énoncez précisément le canevas commun aux trois
démonstrations (PRF-02 comprise) --- ce que l'on suppose, quel genre d'impossibilité on
produit, et quel fait arithmétique sur les entiers fait tourner le moteur.
\item[(d)] Faites tourner le canevas sur $ \sqrt{4} $ et désignez la ligne exacte où il casse.
S'il ne cassait pas, quelque chose irait très mal --- pourquoi ?
\end{enumerate}

\end{problem}

\noindent Une démonstration qu'on ne sait exécuter qu'une fois est un tour de main ;
une démonstration qu'on sait exécuter sur toute une famille est une méthode. Extraire le
canevas réutilisable d'un argument déjà rédigé --- et le mettre à l'épreuve sur un cas où la
conclusion est fausse --- voilà comment les mathématiciens apprennent réellement les
techniques. La question (b) est une favorite parce que le canevas se transplante sur un
énoncé de logarithmes où nulle racine carrée n'est en vue : l'impossibilité y devient un
conflit de parités entre puissances de 2 et puissances de 3.

\begin{refs}
\item Contexte : Raisonnement par l'absurde --- Wikipédia (\url{https://fr.wikipedia.org/wiki/Raisonnement_par_l%27absurde})
\item Lecture : R. Hammack, \textit{Book of Proof}, ch. 6
\item Cours : MIT OCW 6.042J Mathematics for Computer Science (en anglais) (\url{https://ocw.mit.edu/courses/6-042j-mathematics-for-computer-science-fall-2010/})
\end{refs}

\bigskip

\begin{problem}[{Récurrence : les n premiers nombres impairs --- Lycée}]
\textbf{PRF-04.} \textbf{Problème.} Pour tout entier strictement positif n,
\[ 1 + 3 + 5 + \cdots + (2n-1) = n^2 . \]
\begin{enumerate}
\item[(a)] Démontrez l'identité par récurrence, en énonçant intégralement l'initialisation et
l'hérédité.
\item[(b)] Trouvez ensuite une seconde démonstration qui n'utilise aucune récurrence --- une façon
de \textit{voir} l'identité --- et comparez ce que chacune des deux exige du lecteur.
\end{enumerate}

\end{problem}

\noindent L'identité elle-même est très ancienne : les pythagoriciens disposaient des
cailloux en équerres (gnomons) autour d'un carré et le regardaient croître. La récurrence
comme principe de démonstration explicitement énoncé remonte d'ordinaire à Francesco
Maurolico (1575) --- qui a démontré cette identité même --- et à Pascal ; le nom \og{} induction
mathématique \fg{} a été fixé par Augustus De Morgan en 1838. La récurrence est la manière
standard de gravir une échelle infinie barreau par barreau, et voici son premier exercice
traditionnel.

\begin{refs}
\item Contexte : Raisonnement par récurrence --- Wikipédia (\url{https://fr.wikipedia.org/wiki/Raisonnement_par_r%C3%A9currence})
\item Lecture : R. Hammack, \textit{Book of Proof}, ch. 10
\item Cours : MIT OCW 6.042J Mathematics for Computer Science (en anglais) (\url{https://ocw.mit.edu/courses/6-042j-mathematics-for-computer-science-fall-2010/})
\end{refs}

\bigskip

\begin{problem}[{Trouvez la faille : tous les chevaux sont de la même couleur --- Lycée}]
\textbf{PRF-05.} \textbf{Problème.} Voici une \og{} démonstration \fg{} par récurrence du fait que tous les
chevaux sont de la même couleur. \textit{Initialisation : tout ensemble de 1 cheval ne contient
que des chevaux d'une seule couleur. Hérédité : supposons que tout ensemble de n chevaux
soit monochrome. Étant donné n+1 chevaux, mettez-en un de côté ; les n restants partagent
une couleur. Remettez-le et mettez de côté un cheval différent ; de nouveau les n restants
partagent une couleur. Les deux groupes se recouvrent, donc les n+1 chevaux partagent une
seule et même couleur. Donc, par récurrence, tous les chevaux sont de la même couleur.}
La conclusion est absurde, donc la démonstration est fausse. Votre tâche n'est pas de le
constater, mais de localiser la défaillance.
\begin{enumerate}
\item[(a)] Identifiez la valeur exacte de n à laquelle l'argument d'hérédité casse pour la
première fois, et dites précisément quelle phrase y échoue, et pourquoi.
\item[(b)] Expliquez pourquoi \og{} mais les chevaux ont des couleurs variées \fg{} n'est pas une réponse
recevable.
\end{enumerate}

\end{problem}

\noindent Cette parabole a été popularisée par George Pólya dans les années 1950 (sa
version d'origine démontrait que toutes les jeunes filles ont les yeux de la même couleur).
C'est le vaccin classique contre une maladie précise : une hérédité valable pour les grands
n mais qui suppose en silence quelque chose de faux pour un unique petit n. Déboguer une
démonstration est un savoir-faire distinct de celui de l'écrire, et il s'apprend le mieux sur
des arguments conçus pour échouer.

\begin{refs}
\item Contexte : Tous les chevaux sont de la même couleur --- Wikipédia (en anglais) (\url{https://en.wikipedia.org/wiki/All_horses_are_the_same_color})
\item Lecture : G. Pólya, \textit{Induction and Analogy in Mathematics} (1954)
\end{refs}

\bigskip

\begin{problem}[{Trouvez la faille : une démonstration que 1 = 2 --- Lycée}]
\textbf{PRF-06.} \textbf{Problème.} Considérez la \og{} démonstration \fg{} classique suivante. \textit{Soient
$ a = b $ non nuls. Alors $ a^2 = ab $, donc $ a^2 - b^2 = ab - b^2 $, donc
$ (a-b)(a+b) = b(a-b) $, donc $ a + b = b $, donc $ 2b = b $, et par conséquent
$ 2 = 1 $.}
\begin{enumerate}
\item[(a)] Identifiez l'étape exacte qui est illicite et énoncez la règle générale qu'elle
viole.
\item[(b)] Énoncez la leçon sous forme d'un principe précis sur la manipulation des équations :
quelles opérations préservent la vérité sans condition, et lesquelles portent des
conditions annexes ?
\item[(c)] Construisez votre propre démonstration fallacieuse --- d'un énoncé faux de votre choix ---
qui dissimule le même coup illicite au moins deux étapes plus profond que celle-ci.
\end{enumerate}

\end{problem}

\noindent Les sophismes algébriques de ce genre circulent comme folklore scolaire
depuis le dix-neuvième siècle ; E. A. Maxwell en a rassemblé les classiques dans
\textit{Fallacies in Mathematics} (1959). Ce sont plus que des plaisanteries : chaque ligne
d'une vraie démonstration est une affirmation assortie d'une justification, et le plus
rapide moyen d'intérioriser cette discipline est de regarder une chaîne d'équations à l'air
assuré y faire entrer un mensonge en fraude.

\begin{refs}
\item Contexte : Pseudo-démonstration d'égalité entre nombres --- Wikipédia (\url{https://fr.wikipedia.org/wiki/Pseudo-d%C3%A9monstration_d%27%C3%A9galit%C3%A9_entre_nombres})
\item Lecture : E. A. Maxwell, \textit{Fallacies in Mathematics} (Cambridge, 1959)
\end{refs}

\bigskip

\begin{problem}[{Principe des tiroirs : cinq points entiers --- Lycée}]
\textbf{PRF-07.} \textbf{Problème.} Un point entier du plan est un point dont les coordonnées sont
entières.
\begin{enumerate}
\item[(a)] Démontrez que parmi 5 points entiers quelconques, il en existe deux dont le milieu est
lui-même un point entier.
\item[(b)] Montrez que 5 ne peut pas être remplacé par 4.
\item[(c)] Déterminez ensuite, avec démonstration, le plus petit nombre de points entiers de
l'espace de dimension trois qui force l'existence d'une telle paire.
\end{enumerate}

\end{problem}

\noindent Le principe des tiroirs --- si l'on range plus de k objets dans k tiroirs,
l'un des tiroirs en contient deux --- a l'air trop évident pour être une arme. Dirichlet a
montré le contraire, en se servant de son \textit{Schubfachprinzip} (\og{} principe des tiroirs \fg{},
1834) pour démontrer des faits profonds sur l'approximation des irrationnels par des
fractions. Tout l'art réside dans le choix des tiroirs ; ce classique du point milieu,
standard des olympiades depuis des décennies, en est l'illustration la plus limpide
possible.

\begin{refs}
\item Contexte : Principe des tiroirs --- Wikipédia (\url{https://fr.wikipedia.org/wiki/Principe_des_tiroirs})
\item Lecture : A. Engel, \textit{Problem-Solving Strategies}, ch. 4 (le principe des tiroirs)
\end{refs}

\bigskip

\begin{problem}[{Double dénombrement : le lemme des poignées de main --- Lycée}]
\textbf{PRF-08.} \textbf{Problème.} Démontrez que, dans toute assemblée, le nombre de
personnes ayant serré un nombre impair de mains est pair. (De façon équivalente : dans tout
graphe fini, la somme des degrés des sommets vaut deux fois le nombre d'arêtes, de sorte que
le nombre de sommets de degré impair est pair.) Arme suggérée : comptez une même quantité ---
les participations à une poignée de main --- de deux manières différentes.
\end{problem}

\noindent C'est le premier théorème de la théorie des graphes, et il figure dans le
premier article de la théorie des graphes : la résolution par Euler, en 1736, du problème des
ponts de Königsberg. La technique, le double dénombrement, est l'un des gestes les plus
élégants de la combinatoire --- une égalité démontrée non par manipulation, mais en décrivant
un même ensemble depuis deux points de vue. Si simple soit-il, le lemme des poignées de main
tue quantité de constructions à l'air plausible (essayez de dessiner un graphe ayant
exactement trois sommets de degré 3 et aucun autre).

\begin{refs}
\item Contexte : Lemme des poignées de main --- Wikipédia (\url{https://fr.wikipedia.org/wiki/Lemme_des_poign%C3%A9es_de_main})
\item Contexte : Preuve par double dénombrement --- Wikipédia (\url{https://fr.wikipedia.org/wiki/Preuve_par_double_d%C3%A9nombrement})
\item Cours : MIT OCW 6.042J Mathematics for Computer Science (en anglais) (\url{https://ocw.mit.edu/courses/6-042j-mathematics-for-computer-science-fall-2010/})
\end{refs}

\bigskip

\begin{problem}[{Invariants : l'échiquier mutilé --- Lycée}]
\textbf{PRF-09.} \textbf{Problème.} Retirez deux cases d'angle diagonalement opposées d'un échiquier
ordinaire 8 $\times$ 8, ce qui laisse 62 cases. Un domino recouvre exactement deux cases partageant
un côté.
\begin{enumerate}
\item[(a)] Démontrez que l'échiquier restant ne peut pas être pavé par 31 dominos.
\item[(b)] Tranchez ensuite la question réciproque : si l'on retire à la place une case noire et
une case blanche --- n'importe laquelle de ces paires --- les 62 cases restantes peuvent-elles
toujours être pavées ? Déterminez la réponse, avec démonstration.
\end{enumerate}

\end{problem}

\noindent Posé par le philosophe Max Black dans \textit{Critical Thinking} (1946) et
chéri depuis lors ; Gomory a répondu à la seconde question par un argument resté célèbre. La
méthode --- trouver une quantité ou une configuration que tout coup licite préserve, puis
montrer que le départ et l'arrivée n'en donnent pas la même valeur --- c'est le principe
d'invariance, une bête de somme qui va des mathématiques récréatives jusqu'aux lois de
conservation. John McCarthy a proposé l'échiquier mutilé à la jeune communauté de
l'intelligence artificielle en 1964 comme cas d'école, précisément parce que la recherche
exhaustive y est immense alors qu'une seule idée clôt l'affaire à l'instant.

\begin{refs}
\item Contexte : Problème de l'échiquier mutilé --- Wikipédia (\url{https://fr.wikipedia.org/wiki/Probl%C3%A8me_de_l%27%C3%A9chiquier_mutil%C3%A9})
\item Lecture : A. Engel, \textit{Problem-Solving Strategies}, ch. 1--2
\end{refs}

\bigskip

\begin{problem}[{Ce que la démonstration par l'exemple ne peut pas faire --- Lycée}]
\textbf{PRF-10.} \textbf{Problème.} Trois exercices sur ce que la vérification d'exemples peut et ne peut
pas établir.
\begin{enumerate}
\item[(a)] Énoncez précisément ce qui cloche dans la \og{} démonstration par l'exemple \fg{} : pourquoi
aucune liste finie de cas vérifiés ne peut-elle établir un énoncé de la forme \og{} pour tout
n \fg{} ? Dites aussi ce que la vérification d'exemples \textit{peut} légitimement démontrer.
\item[(b)] Le polynôme $ n^2 + n + 41 $ produit des nombres premiers pour
$ n = 0, 1, 2, \dots, 39 $ --- quarante d'affilée. Trouvez, avec démonstration, le plus
petit n pour lequel il échoue, et expliquez pourquoi il était condamné à échouer, sans rien
calculer.
\item[(c)] Fermat croyait que tout nombre $ F_k = 2^{2^k} + 1 $ est premier, et les cinq
premiers le sont. Montrez que $ 641 \mid F_5 $ en utilisant les identités
$ 641 = 5^4 + 2^4 = 5 \cdot 2^7 + 1 $ plutôt qu'une division posée.
\end{enumerate}

\end{problem}

\noindent Les exemples sont ce qui fait naître les conjectures, et les exemples sont
ce qui a tué celle-ci : en 1732 le jeune Euler a factorisé $ F_5 = 4294967297 $, démolissant
la conjecture à laquelle Fermat s'était engagé avec le plus d'assurance. (Le polynôme de la
question (b) est lui aussi d'Euler, et date de 1772.) \textit{Mathematics and Plausible
Reasoning} de Pólya est la méditation classique sur le juste rapport entre indices et
démonstration : les exemples suggèrent ; seul l'argument établit.

\begin{refs}
\item Contexte : Nombre de Fermat --- Wikipédia (\url{https://fr.wikipedia.org/wiki/Nombre_de_Fermat})
\item Lecture : G. Pólya, \textit{Mathematics and Plausible Reasoning} (1954)
\end{refs}

\bigskip

\begin{problem}[{Quand une figure est-elle une démonstration ? --- Lycée}]
\textbf{PRF-11.} \textbf{Problème.} La figure classique de réarrangement pour le théorème de Pythagore
place quatre copies d'un triangle rectangle de côtés de l'angle droit a et b et d'hypoténuse
c à l'intérieur d'un carré de côté $ a + b $, de deux manières différentes. L'une laisse à
découvert un unique carré incliné d'aire $ c^2 $ ; l'autre laisse deux carrés d'aires
$ a^2 $ et $ b^2 $.
\begin{enumerate}
\item[(a)] Dessinez les deux dispositions et rédigez, en phrases complètes, l'argument que les
figures encodent --- y compris chaque fait que la figure vous demande d'admettre sur parole
(que les pièces ne se chevauchent pas, que les régions découvertes sont bien des carrés,
que l'aire survit au réarrangement).
\item[(b)] Proposez des critères : quand un dessin doit-il compter comme une démonstration, et
quand n'est-il qu'un indice ?
\item[(c)] Le célèbre paradoxe de découpage dans lequel un carré 8 $\times$ 8 est coupé en quatre pièces
puis réassemblé en un rectangle 5 $\times$ 13 \og{} démontre \fg{} que 64 = 65. Expliquez exactement où
cette figure ment, et ce que vos critères de la question (b) en disent.
\end{enumerate}

\end{problem}

\noindent Les démonstrations du théorème de Pythagore par réarrangement sont très
anciennes --- l'une accompagne le \textit{Zhoubi Suanjing} en Chine, et l'on rapporte que
Bhāskara II, dans l'Inde du douzième siècle, dessina la figure et écrivit un seul mot :
\og{} Vois ! \fg{}. Les mathématiciens débattent depuis lors de la confiance que l'on peut accorder à
une figure ; la réponse honnête est qu'un bon dessin est un argument comprimé dont la
décompression incombe au lecteur. \textit{Proofs Without Words} de Nelsen en rassemble une
centaine de beaux spécimens sur lesquels s'exercer.

\begin{refs}
\item Contexte : Preuve sans mots --- Wikipédia (\url{https://fr.wikipedia.org/wiki/Preuve_sans_mots})
\item Contexte : Théorème de Pythagore --- Wikipédia (\url{https://fr.wikipedia.org/wiki/Th%C3%A9or%C3%A8me_de_Pythagore})
\item Lecture : R. B. Nelsen, \textit{Proofs Without Words: Exercises in Visual Thinking} (MAA, 1993)
\end{refs}

\bigskip

\begin{problem}[{La symétrie comme stratégie : des pièces sur une table ronde --- Lycée}]
\textbf{PRF-37.} \textbf{Problème.} Deux joueurs posent tour à tour des pièces circulaires identiques sur
une table ronde. Chaque pièce doit reposer entièrement sur la table sans recouvrir aucune
pièce déjà posée (le contact est autorisé), et le joueur qui n'a plus de coup licite
perd.
\begin{enumerate}
\item[(a)] Démontrez que le premier joueur peut toujours gagner, quelles que soient les tailles
de la table et des pièces : décrivez une stratégie et démontrez qu'elle ne manque jamais
de coups avant l'adversaire.
\item[(b)] Désignez les faits géométriques exacts qu'utilise votre démonstration --- où,
précisément, l'argument a-t-il besoin de la symétrie de la table, et pourquoi le premier
coup doit-il être le centre ?
\item[(c)] Pour une table carrée, une table rectangulaire non carrée et une table en L, décidez
si votre démonstration survit ; partout où elle échoue, identifiez l'étape exacte qui
meurt.
\end{enumerate}

\end{problem}

\noindent Un classique des mathématiques récréatives qui distille un principe
sérieux : une symétrie du plateau peut être enrôlée comme stratégie, car l'image miroir de
tout coup licite de l'adversaire est licite à coup sûr --- pourvu que l'on se soit d'abord
emparé de l'unique point que la symétrie laisse fixe. Les stratégies d'appariement et de
miroir de ce genre décident quantité de jeux à deux joueurs sans le moindre calcul, et cette
même habitude de pensée --- argumenter une fois, laisser la symétrie répéter l'argument à
votre place --- raccourcit les démonstrations partout en mathématiques, depuis chaque \og{} sans
perte de généralité \fg{} en géométrie jusqu'au dénombrement par involution de PRF-47.

\begin{refs}
\item Contexte : Argument de vol de stratégie --- Wikipédia (en anglais) (\url{https://en.wikipedia.org/wiki/Strategy-stealing_argument})
\item Contexte : La symétrie en mathématiques --- Wikipédia (en anglais) (\url{https://en.wikipedia.org/wiki/Symmetry_in_mathematics})
\item Lecture : E. R. Berlekamp, J. H. Conway \& R. K. Guy, \textit{Winning Ways for Your Mathematical Plays} (1982)
\item Lecture : A. Engel, \textit{Problem-Solving Strategies}, ch. 13 (les jeux)
\end{refs}

\bigskip

\begin{problem}[{Démontrez ou réfutez : irrationnel + irrationnel --- Lycée, short}]
\textbf{PRF-38.} \textbf{Problème.} Démontrez ou réfutez : la somme de deux nombres
irrationnels est irrationnelle. Faites ensuite de même pour le produit de deux nombres
irrationnels.
\end{problem}

\noindent \og{} Démontrez ou réfutez \fg{} est un genre à part entière : la moitié de la
tâche consiste à décider de quel côté on se range, et une réfutation est une démonstration
elle aussi --- un contre-exemple bien choisi, présenté avec le soin dû à un théorème (PRF-43
en est l'atelier complet). Remarquez à quel point les deux charges sont inégales : une
affirmation universelle meurt d'un seul exemple.

\begin{refs}
\item Contexte : Nombre irrationnel --- Wikipédia (\url{https://fr.wikipedia.org/wiki/Nombre_irrationnel})
\item Lecture : R. Hammack, \textit{Book of Proof}, ch. 9 (la réfutation)
\end{refs}

\bigskip

\begin{problem}[{Trois démonstrations que 6 divise n$^3$ $-$ n --- Lycée, short}]
\textbf{PRF-39.} \textbf{Problème.} Démontrez que $ n^3 - n $ est divisible par 6 pour
tout entier n --- trois fois, par trois arguments véritablement différents (candidats : la
récurrence ; la factorisation en entiers consécutifs ; l'arithmétique modulo 2 et modulo
3).
\end{problem}

\noindent Une miniature de la discipline récurrente de cette page (PRF-15, PRF-18,
PRF-24) : une vérité, plusieurs démonstrations, et une comparaison de ce que chacune voit.
L'argument par entiers consécutifs se généralise aux coefficients binomiaux ; l'argument
modulaire est le germe du petit théorème de Fermat, qui fait pour $ n^p - n $ ce que ce
problème fait pour $ n^3 - n $ (PRF-18).

\begin{refs}
\item Contexte : Arithmétique modulaire --- Wikipédia (\url{https://fr.wikipedia.org/wiki/Arithm%C3%A9tique_modulaire})
\item Lecture : R. Hammack, \textit{Book of Proof}, ch. 10
\end{refs}

\bigskip

\begin{problem}[{Six personnes à une soirée, en dix phrases --- Lycée, short}]
\textbf{PRF-40.} \textbf{Problème.} Démontrez que parmi six personnes quelconques, il en
existe trois qui se connaissent deux à deux, ou trois qui sont deux à deux étrangères l'une
à l'autre --- en dix phrases au plus. Montrez ensuite que six est optimal : disposez cinq
personnes de sorte qu'aucun des deux trios n'existe.
\end{problem}

\noindent Le théorème des amis et des étrangers --- R(3,3) = 6 dans la notation de
Ramsey --- figurait au concours Putnam de 1953 et ouvre toute visite guidée de la théorie de
Ramsey, l'étude de l'ordre qu'aucune quantité de désordre ne peut éviter. Le budget de
phrases est ici le véritable exercice : un argument correct s'étale facilement, et le
comprimer sans rien perdre, c'est l'écriture de démonstration sous sa forme la plus pure. La
frontière est étonnamment proche : la valeur de R(5,5) est toujours inconnue.

\begin{refs}
\item Contexte : Théorème des amis et des étrangers --- Wikipédia (\url{https://fr.wikipedia.org/wiki/Th%C3%A9or%C3%A8me_des_amis_et_des_%C3%A9trangers})
\item Contexte : Théorème de Ramsey --- Wikipédia (\url{https://fr.wikipedia.org/wiki/Th%C3%A9or%C3%A8me_de_Ramsey})
\end{refs}

\bigskip

\begin{problem}[{La réciproque n'est pas gratuite --- Lycée}]
\textbf{PRF-46.} \textbf{Problème.} Une affirmation de la forme \og{} P si et seulement si Q \fg{} énonce deux
implications distinctes --- $ P \Rightarrow Q $ et sa réciproque $ Q \Rightarrow P $ --- et
aucune des deux ne découle de l'autre.
\begin{enumerate}
\item[(a)] Démontrez qu'un entier n est divisible par 3 si et seulement si la somme de ses
chiffres décimaux est divisible par 3. Rédigez les deux sens comme deux arguments
clairement étiquetés, et marquez à quel sens appartient chaque étape de votre travail.
\item[(b)] Pour chacun des énoncés ci-dessous, écrivez sa réciproque et décidez, par une
démonstration ou par un contre-exemple, si cette réciproque est vraie : (i) si un
quadrilatère est un carré, alors ses diagonales sont perpendiculaires ; (ii) si n est un
nombre premier supérieur à 2, alors n est impair ; (iii) si un entier est divisible par 24,
alors il est divisible à la fois par 4 et par 6.
\item[(c)] Certaines équivalences se démontrent d'un seul trait, par une chaîne d'étapes toutes
réversibles. Déterminez toutes les solutions réelles de
\[ \sqrt{x+7} \;=\; x+1, \]
puis repérez l'étape exacte de la méthode habituelle \og{} on élève les deux membres au
carré \fg{} qui n'est \textit{pas} réversible, et expliquez pourquoi cette seule étape oblige la
méthode à s'achever par une vérification.
\end{enumerate}

\end{problem}

\noindent Euclide traite déjà un théorème et sa réciproque comme deux travaux
distincts : la proposition I.5 des \textit{Éléments} démontre qu'un triangle isocèle a des
angles à la base égaux, et I.6 --- proposition entièrement séparée, munie de sa propre
démonstration --- en établit la réciproque. La plupart des erreurs d'étudiants du genre \og{} je
l'ai démontré à l'envers \fg{} sont des erreurs de réciproque, et le phénomène des racines
parasites de la question (c) est la même faute en habits algébriques : élever au carré est
une voie à sens unique, de sorte que la chaîne d'équations démontre seulement que \textit{si}
une solution existe, elle figure parmi les candidates trouvées. L'abréviation \og{} iff \fg{} est de
Paul Halmos, et la convention selon laquelle le mot \og{} si \fg{} dans une \textit{définition}
signifie tacitement \og{} si et seulement si \fg{} est de celles que tout lecteur doit assimiler
tôt.

\begin{refs}
\item Contexte : Condition nécessaire et suffisante --- Wikipédia (\url{https://fr.wikipedia.org/wiki/Condition_n%C3%A9cessaire_et_suffisante})
\item Contexte : Implication réciproque --- Wikipédia (\url{https://fr.wikipedia.org/wiki/Implication_r%C3%A9ciproque})
\item Lecture : R. Hammack, \textit{Book of Proof}, le chapitre sur la démonstration des énoncés non conditionnels (gratuit en ligne (\url{https://richardhammack.github.io/BookOfProof/}))
\item Lecture : D. J. Velleman, \textit{How to Prove It}, ch. 3
\end{refs}

\bigskip

\begin{problem}[{Deux quantificateurs, deux sens --- Lycée, short}]
\textbf{PRF-54.} \textbf{Problème.} Décidez, avec démonstration, si les deux phrases de
chacune des paires ci-dessous disent la même chose : (i) \og{} pour tout entier x il existe un
entier y tel que $ y > x $ \fg{} contre \og{} il existe un entier y tel que $ y > x $ pour tout
entier x \fg{} ; (ii) \og{} tout réel strictement positif est plus petit qu'un certain entier \fg{}
contre \og{} un certain entier est plus grand que tout réel strictement positif \fg{}. Énoncez
ensuite, en une phrase, la règle générale sur l'échange de \og{} pour tout \fg{} et \og{} il existe \fg{}, et
identifiez le seul sens de cet échange qui soit toujours valide.
\end{problem}

\noindent L'ordre des quantificateurs est de loin la première source de
démonstrations fausses qui ont l'air justes, et l'ennui est que la langue courante le
masque : \og{} tout le monde aime quelqu'un \fg{} et \og{} quelqu'un est aimé de tout le monde \fg{} ne
diffèrent que par l'accent lorsqu'on les dit, et totalement par le sens lorsqu'on les écrit.
Le \textit{Begriffsschrift} de Frege (1879) fut la première notation capable d'exprimer sans
ambiguïté une quantification imbriquée, et ce n'est pas un hasard si l'analyse rigoureuse ---
où continuité, continuité uniforme et convergence ne se distinguent les unes des autres que
par l'ordre des quantificateurs --- est devenue possible au même siècle. PRF-44 met le même
savoir-faire au travail sur la négation d'un énoncé en $\varepsilon$--$\delta$.

\begin{refs}
\item Contexte : Calcul des prédicats --- Wikipédia (\url{https://fr.wikipedia.org/wiki/Calcul_des_pr%C3%A9dicats})
\item Lecture : D. J. Velleman, \textit{How to Prove It}, ch. 2
\item Lecture : R. Hammack, \textit{Book of Proof}, ch. 2 (gratuit en ligne (\url{https://richardhammack.github.io/BookOfProof/}))
\end{refs}

\bigskip

\begin{problem}[{Le principe des tiroirs déguisé : inventer les tiroirs --- Lycée}]
\textbf{PRF-55.} \textbf{Problème.} Le principe des tiroirs est trivial à énoncer et difficile à employer,
car toute la difficulté est d'inventer les tiroirs. Dans PRF-07 les tiroirs étaient presque
visibles ; dans chacune des questions ci-dessous, ils ne vous sont pas donnés du tout.
\begin{enumerate}
\item[(a)] Démontrez que parmi $ n+1 $ nombres quelconques choisis dans
$ \{1, 2, \dots, 2n\} $, l'un d'eux en divise un autre.
\item[(b)] Démontrez que toute liste de n entiers (pas nécessairement distincts) contient un bloc
non vide de termes consécutifs dont la somme est divisible par n.
\item[(c)] Démontrez le théorème d'Erdős--Szekeres : toute suite de $ n^2 + 1 $ réels distincts
contient une sous-suite croissante de longueur $ n+1 $ ou une sous-suite décroissante de
longueur $ n+1 $.
\item[(d)] Démontrez le théorème d'approximation de Dirichlet : pour tout réel $\alpha$ et tout entier
strictement positif N, il existe des entiers p et q avec $ 1 \le q \le N $ et
\[ \left| \alpha - \frac{p}{q} \right| < \frac{1}{qN}. \]
\item[(e)] Voici maintenant le but de l'exercice. Pour chacune des questions (a)--(d), écrivez
explicitement quels étaient les objets à ranger, quels étaient les tiroirs, et combien il y
en avait de chaque --- puis dites, en une phrase par question, ce qui rendait ces tiroirs
difficiles à voir.
\end{enumerate}

\end{problem}

\noindent Dirichlet a introduit le \textit{Schubfachprinzip} en 1834 non comme une
curiosité mais comme un outil, et la question (d) est ce pour quoi il l'a bâti : un énoncé
absolument non évident sur l'approximation des irrationnels, démontré en laissant tomber
$ N+1 $ parties fractionnaires dans N tiroirs. La question (c) est le théorème d'Erdős et
Szekeres de 1935, dont la démonstration la plus habile --- d'ordinaire attribuée à Seidenberg
(1959) --- étiquette chaque terme par un couple de nombres et laisse le principe des tiroirs
conclure en une ligne. Le schéma est le même dans tous les cas, et il mérite d'être nommé :
les tiroirs ne sont pas les objets qu'on vous a remis, mais un invariant calculé à partir
d'eux --- une partie impaire, un reste, un couple de longueurs, une partie fractionnaire.

\begin{refs}
\item Contexte : Théorème d'Erdős-Szekeres --- Wikipédia (\url{https://fr.wikipedia.org/wiki/Th%C3%A9or%C3%A8me_d%27Erd%C5%91s-Szekeres})
\item Contexte : Approximation diophantienne --- Wikipédia (\url{https://fr.wikipedia.org/wiki/Approximation_diophantienne})
\item Lecture : A. Engel, \textit{Problem-Solving Strategies}, ch. 4 (le principe des tiroirs)
\item Lecture : M. Aigner \& G. M. Ziegler, \textit{Proofs from THE BOOK}, le chapitre sur le principe des tiroirs et le double dénombrement
\end{refs}

\bigskip


\section*{Licence}

\begin{problem}[{Récurrence forte : l'existence de la factorisation en nombres premiers --- Licence}]
\textbf{PRF-12.} \textbf{Problème.} Tout entier $ n \ge 2 $ peut s'écrire comme un produit de nombres
premiers.
\begin{enumerate}
\item[(a)] Démontrez-le par récurrence forte.
\item[(b)] Démontrez ensuite le même énoncé d'une seconde manière : supposez qu'il soit en
défaut, considérez le \textit{plus petit} contre-exemple, et aboutissez à une
contradiction.
\item[(c)] Enfin, expliquez pourquoi ces deux démonstrations sont la même démonstration sous des
habits différents --- autrement dit, explicitez le rapport entre la récurrence forte et le
principe du bon ordre (toute partie non vide d'entiers strictement positifs possède un plus
petit élément).
\end{enumerate}

\end{problem}

\noindent L'existence de la factorisation en nombres premiers est la moitié facile
du théorème fondamental de l'arithmétique ; l'unicité est plus subtile et a dû attendre les
\textit{Disquisitiones Arithmeticae} de Gauss (1801) pour être énoncée et démontrée en toute
généralité. Le point pédagogique est ici l'interchangeabilité de la récurrence forte et du
plus petit contre-exemple --- deux tournures pour un seul principe --- et l'apprentissage d'une
écriture aisée des deux. L'unicité, par le lemme d'Euclide, est une suite qui en vaut la
peine : voyez la lecture.

\begin{refs}
\item Contexte : Théorème fondamental de l'arithmétique --- Wikipédia (\url{https://fr.wikipedia.org/wiki/Th%C3%A9or%C3%A8me_fondamental_de_l%27arithm%C3%A9tique})
\item Lecture : G. H. Hardy \& E. M. Wright, \textit{An Introduction to the Theory of Numbers}, ch. 1--2
\item Cours : MIT OCW 6.042J Mathematics for Computer Science (en anglais) (\url{https://ocw.mit.edu/courses/6-042j-mathematics-for-computer-science-fall-2010/})
\end{refs}

\bigskip

\begin{problem}[{Le principe extrémal --- Licence}]
\textbf{PRF-13.} \textbf{Problème.} Le principe extrémal : pour démontrer qu'il existe un objet possédant
une certaine propriété, examinez un objet extrême --- le plus grand, le plus petit, le plus
long, le plus proche --- et montrez que s'il n'avait pas cette propriété, on pourrait le
rendre plus extrême encore, ce qui est contradictoire. Employez-le deux fois.
\begin{enumerate}
\item[(a)] Un tournoi est un graphe complet dont chaque arête a reçu une orientation (chacune des
paires parmi n joueurs s'est rencontrée une fois ; quelqu'un a gagné chaque partie).
Démontrez que tout tournoi fini contient un chemin hamiltonien : un ordre sur l'ensemble
des joueurs dans lequel chacun a battu le suivant. Considérez un chemin orienté de longueur
maximale.
\item[(b)] Étant donné n points rouges et n points bleus du plan, jamais trois alignés, démontrez
qu'on peut les apparier, chaque rouge avec un bleu, de sorte que les n segments de liaison
soient deux à deux sans intersection.
\end{enumerate}

\end{problem}

\noindent La question (a) est un théorème de László Rédei (1934). Le principe
extrémal, c'est le principe du bon ordre en tenue de travail, et il est d'une polyvalence
étonnante : la célèbre démonstration par Kelly du théorème de Sylvester--Gallai (PRF-26) est
un argument extrémal pur. Le chapitre qu'Engel lui consacre est le terrain d'entraînement
classique.

\begin{refs}
\item Contexte : Tournoi (théorie des graphes) --- Wikipédia (\url{https://fr.wikipedia.org/wiki/Tournoi_(th%C3%A9orie_des_graphes)})
\item Contexte : Principe du bon ordre --- Wikipédia (en anglais) (\url{https://en.wikipedia.org/wiki/Well-ordering_principle})
\item Lecture : A. Engel, \textit{Problem-Solving Strategies}, ch. 3 (le principe extrémal)
\end{refs}

\bigskip

\begin{problem}[{Euclide : les nombres premiers sont en infinité --- Licence}]
\textbf{PRF-14.} \textbf{Problème.} Le théorème d'Euclide --- et l'art de ne pas le citer de travers.
\begin{enumerate}
\item[(a)] Démontrez qu'il existe une infinité de nombres premiers.
\item[(b)] Une lecture erronée très répandue de l'argument d'Euclide prétend que si
$ p_1, \dots, p_n $ sont les n premiers nombres premiers, alors
$ p_1 p_2 \cdots p_n + 1 $ est lui-même premier. Réfutez-la : trouvez, avec démonstration,
le plus petit n pour lequel $ p_1 p_2 \cdots p_n + 1 $ est composé.
\item[(c)] Réécrivez votre démonstration de (a) de sorte qu'elle n'affirme rien que votre
réfutation de (b) contredise --- énoncez exactement ce que le nombre auxiliaire de l'argument
est garanti de faire, et rien de plus.
\end{enumerate}

\end{problem}

\noindent Proposition 20 du livre IX des \textit{Éléments} d'Euclide (vers 300 av.
J.-C.) : \og{} les nombres premiers sont en quantité plus grande que toute multitude assignée de
nombres premiers \fg{} --- pour bien des mathématiciens, la première démonstration qu'ils appellent
belle. L'argument propre à Euclide est direct, et non le raisonnement par l'absurde sous
lequel on le raconte si souvent ; ce sont les récits qui engendrent la lecture erronée de la
question (b). Michael Hardy et Catherine Woodgold ont documenté ce téléphone arabe dans
\og{} Prime Simplicity \fg{}. Restituer Euclide exactement est un rite de passage de l'écriture
soignée des démonstrations --- et le point de départ de deux démonstrations plus profondes du
même théorème dans la section Master (PRF-28, PRF-29).

\begin{refs}
\item Contexte : Théorème d'Euclide sur les nombres premiers --- Wikipédia (\url{https://fr.wikipedia.org/wiki/Th%C3%A9or%C3%A8me_d%27Euclide_sur_les_nombres_premiers})
\item Lecture : M. Aigner \& G. M. Ziegler, \textit{Proofs from THE BOOK}, ch. 1
\item Article : M. Hardy \& C. Woodgold, \og{} Prime Simplicity \fg{} (2009), \textit{Mathematical Intelligencer} 31
\end{refs}

\bigskip

\begin{problem}[{Le problème de Bâle --- Licence}]
\textbf{PRF-15.} \textbf{Problème.} Considérez la série $ \sum_{n=1}^{\infty} \frac{1}{n^2} $ ---
la pièce maîtresse de cette page.
\begin{enumerate}
\item[(a)] Montrez que la série converge et que sa valeur est inférieure à 2.
\item[(b)] Le problème de Bâle : démontrez que
\[ \sum_{n=1}^{\infty} \frac{1}{n^2} \;=\; \frac{\pi^2}{6}. \]
Toute voie honnête est recevable ; deux d'entre elles sont cartographiées dans PRF-16 et
PRF-17. On connaît au moins quatorze démonstrations essentiellement différentes --- après la
première, allez en chercher une seconde dans les références et déterminez précisément où
les deux font un travail différent.
\end{enumerate}

\end{problem}

\noindent La pièce maîtresse de cette page. Pietro Mengoli a posé la question en
1650 : les inverses des carrés se somment manifestement en quelque chose d'un peu supérieur
à 1,6 --- mais en \textit{quoi} ? Elle a vaincu Leibniz et vaincu les Bernoulli ; Jakob
Bernoulli, de Bâle, qui avait sommé tant de séries, publia un appel en 1689 : \og{} si quelqu'un
trouve et nous communique ce qui a jusqu'ici échappé à nos efforts, grande sera notre
gratitude. \fg{} En 1734-35 Leonhard Euler --- né à Bâle, élève de Jean Bernoulli, alors âgé de
vingt-huit ans et installé à Saint-Pétersbourg --- annonça la réponse, et l'apparition
parfaitement inattendue de $\pi$ dans un problème où nul cercle n'était en vue fit son nom dans
toute l'Europe. Le problème a tout ensemencé : Euler évalua ensuite $\zeta$(2k) pour tout k, relia
la fonction zêta aux nombres premiers (PRF-28), et remit ainsi à Riemann l'objet central de
la théorie analytique des nombres. La compilation de Robin Chapman rassemble quatorze
démonstrations ; Aigner et Ziegler en impriment trois des plus belles.

\begin{refs}
\item Contexte : Problème de Bâle --- Wikipédia (\url{https://fr.wikipedia.org/wiki/Probl%C3%A8me_de_B%C3%A2le})
\item Lecture : W. Dunham, \textit{Euler: The Master of Us All} (MAA, 1999)
\item Lecture : M. Aigner \& G. M. Ziegler, \textit{Proofs from THE BOOK}, ch. \og{} Three times $\pi$$^2$/6 \fg{}
\item Article : R. Chapman, \og{} Evaluating $\zeta$(2) \fg{} (1999-2003), quatorze démonstrations, PDF (en anglais) (\url{https://empslocal.ex.ac.uk/people/staff/rjchapma/etc/zeta2.pdf})
\end{refs}

\bigskip

\begin{problem}[{L'heuristique d'Euler : sin x comme polynôme infini --- Licence}]
\textbf{PRF-16.} \textbf{Problème.} L'argument original d'Euler, en 1734, traitait
$ \frac{\sin x}{x} $ comme un polynôme de degré infini. Ses \og{} racines \fg{} sont
$ \pm\pi, \pm 2\pi, \pm 3\pi, \dots $, et donc --- en factorisant sur les racines comme on
le ferait pour un polynôme --- Euler écrivit
\[ \frac{\sin x}{x} \;=\; \prod_{n=1}^{\infty} \left( 1 - \frac{x^2}{n^2 \pi^2} \right). \]
\begin{enumerate}
\item[(a)] En admettant cette identité, comparez le coefficient de $ x^2 $ des deux côtés et
déduisez-en la valeur de $ \sum 1/n^2 $.
\item[(b)] Auditez maintenant l'argument : dressez la liste de chaque étape injustifiée au regard
des normes du dix-huitième siècle. En particulier : exhibez deux fonctions, chacune donnée
par une série entière convergeant sur $\mathbb{R}$ tout entier, ayant exactement les mêmes zéros et
pourtant non égales --- pourquoi cela condamne-t-il la recette naïve \og{} une fonction est le
produit sur ses racines \fg{}, et quelle information supplémentaire un théorème de
factorisation correct doit-il utiliser ?
\item[(c)] Soit rendez l'argument d'Euler rigoureux (la voie moderne est la factorisation de
Weierstrass--Hadamard du sinus), soit vérifiez le contrôle de vraisemblance d'Euler
lui-même : en posant $ x = \pi/2 $ dans le produit, on retrouve la formule de Wallis pour
$\pi$.
\end{enumerate}

\end{problem}

\noindent C'est la lacune la plus instructive de l'histoire de l'analyse. Euler
savait parfaitement que sa factorisation n'était pas démontrée --- c'est pourquoi il n'a cessé
de revenir au problème, publiant d'autres démonstrations en 1741 et 1755 --- mais il savait
aussi que la réponse était juste, l'ayant vérifiée à de nombreuses décimales et confrontée au
produit de Wallis de 1655. L'épisode est l'illustration canonique de ce que Terence Tao
appelle la pensée pré-rigoureuse par opposition à la pensée post-rigoureuse : le bond qui
trouve la vérité, puis la discipline qui l'assure. Weierstrass a fourni le théorème manquant
un siècle et demi plus tard.

\begin{refs}
\item Contexte : Théorème de factorisation de Weierstrass --- Wikipédia (\url{https://fr.wikipedia.org/wiki/Th%C3%A9or%C3%A8me_de_factorisation_de_Weierstrass})
\item Contexte : Produit de Wallis --- Wikipédia (\url{https://fr.wikipedia.org/wiki/Produit_de_Wallis})
\item Lecture : W. Dunham, \textit{Euler: The Master of Us All} (MAA, 1999)
\item Lecture : T. Tao, \og{} There's more to mathematics than rigour and proofs \fg{} (en anglais) (\url{https://terrytao.wordpress.com/career-advice/theres-more-to-mathematics-than-rigour-and-proofs/})
\end{refs}

\bigskip

\begin{problem}[{Une route rigoureuse vers $\pi$$^2$/6 --- Licence}]
\textbf{PRF-17.} \textbf{Problème.} Deux routes pleinement rigoureuses vers la somme de Bâle. Choisissez-en
une et parcourez-la de bout en bout ; pour l'autre, rédigez en un paragraphe une carte du
terrain d'après les références.
\begin{enumerate}
\item[(a)] Route A (intégrale double) : montrez que
\[ \int_0^1 \!\! \int_0^1 \frac{dx\, dy}{1 - xy} \;=\; \sum_{n=1}^{\infty} \frac{1}{n^2}, \]
en justifiant l'interversion de la somme et de l'intégrale, puis évaluez l'intégrale par un
changement de variables bien choisi.
\item[(b)] Route B (séries de Fourier) : calculez les coefficients de Fourier de $ f(x) = x $
sur $ [-\pi, \pi] $ et appliquez l'identité de Parseval.
\end{enumerate}

\end{problem}

\noindent La route A est la démonstration que Tom Apostol a appelée \og{} une
démonstration qui a échappé à Euler \fg{} dans sa note de 1983, avec l'évaluation par changement
de variables étudiée en profondeur par Beukers, Calabi et Kolk (1993), dont la méthode évalue
$\zeta$(2k) pour tout k ; c'est un joyau d'Aigner--Ziegler. La route B est le premier triomphe
classique de l'analyse de Fourier : l'identité de Parseval transforme un calcul
d'orthogonalité en la somme de Bâle, et la même machine livre $\zeta$(4), $\zeta$(6), \dots{} à partir d'autres
fonctions simples. Comparer les deux routes enseigne ce que signifie \og{} des démonstrations
véritablement différentes \fg{} --- elles se généralisent selon des axes entièrement distincts.

\begin{refs}
\item Contexte : Valeurs particulières de la fonction zêta de Riemann --- Wikipédia (\url{https://fr.wikipedia.org/wiki/Valeurs_particuli%C3%A8res_de_la_fonction_z%C3%AAta_de_Riemann})
\item Article : T. M. Apostol, \og{} A proof that Euler missed: evaluating $\zeta$(2) the easy way \fg{} (1983), \textit{Mathematical Intelligencer} 5
\item Article : F. Beukers, E. Calabi \& J. A. C. Kolk, \og{} Sums of generalized harmonic series and volumes \fg{} (1993), \textit{Nieuw Archief voor Wiskunde}
\item Lecture : E. M. Stein \& R. Shakarchi, \textit{Fourier Analysis: An Introduction}
\end{refs}

\bigskip

\begin{problem}[{Le petit théorème de Fermat, de deux façons --- Licence}]
\textbf{PRF-18.} \textbf{Problème.} Soit p un nombre premier. Démontrez que
\[ a^{p-1} \equiv 1 \pmod{p} \]
pour tout entier a non divisible par p --- deux fois.
\begin{enumerate}
\item[(a)] Démonstration 1 (récurrence) : montrez d'abord, par récurrence sur a à l'aide de la
formule du binôme, que $ a^p \equiv a \pmod{p} $ pour tout entier $ a \ge 0 $ ;
déduisez-en le théorème. Il vous faudra savoir quels coefficients binomiaux
$ \binom{p}{k} $ sont divisibles par p --- démontrez-le également.
\item[(b)] Démonstration 2 (théorie des groupes) : montrez que les résidus non nuls modulo p
forment un groupe pour la multiplication, et déduisez le théorème du théorème de Lagrange
sur l'ordre d'un élément.
\item[(c)] Comparez ensuite : que voit chaque démonstration à quoi l'autre est aveugle, et
laquelle se généralise au théorème d'Euler pour un module composé ?
\end{enumerate}

\end{problem}

\noindent Fermat a énoncé le théorème dans une lettre à Frénicle de Bessy en octobre
1640, ajoutant qu'il en enverrait la démonstration \og{} si je ne craignais d'être trop long \fg{}.
Leibniz en possédait une démonstration qu'il n'a jamais publiée ; Euler a imprimé la première
en 1736. Les deux démonstrations demandées ici forment la paire moderne standard, et leur
contraste est tout le propos : la démonstration par récurrence est élémentaire et se suffit à
elle-même, tandis que la démonstration par la théorie des groupes explique \textit{pourquoi} le
théorème est vrai et s'étend instantanément à tout groupe fini. Le petit théorème de Fermat
est aussi le moteur des tests de primalité de PRF-19 et PRF-20 --- et du cryptosystème RSA.

\begin{refs}
\item Contexte : Petit théorème de Fermat --- Wikipédia (\url{https://fr.wikipedia.org/wiki/Petit_th%C3%A9or%C3%A8me_de_Fermat})
\item Lecture : D. M. Burton, \textit{Elementary Number Theory}
\item Cours : MIT OCW 18.781 Theory of Numbers (en anglais) (\url{https://ocw.mit.edu/courses/18-781-theory-of-numbers-spring-2012/})
\end{refs}

\bigskip

\begin{problem}[{Le nombre 341, ou pourquoi le test de Fermat ment --- Licence}]
\textbf{PRF-19.} \textbf{Problème.} Le test de primalité de Fermat : pour tester n, choisissez une base a
avec $ 1 <a <n $ et calculez $ a^{n-1} \bmod n $ ; si le résultat n'est pas 1,
alors n est à coup sûr composé. La réciproque rendrait le test de primalité trivial --- mais
elle est fausse.
\begin{enumerate}
\item[(a)] Démontrez que $ 2^{340} \equiv 1 \pmod{341} $ bien que
$ 341 = 11 \cdot 31 $ soit composé : le plus petit pseudo-premier de base 2.
\item[(b)] Montrez que
le test en base 3, lui, démasque 341 : $ 3^{340} \not\equiv 1 \pmod{341} $.
\item[(c)] Pire
encore est vrai : il existe des nombres composés qui passent le test de Fermat pour
\textit{toute} base première avec eux --- le plus petit est 561. Démontrez que 561 possède cette
propriété, et extrayez de votre démonstration un critère net indiquant quand un composé n
trompe toutes les bases premières avec lui.
\end{enumerate}

\end{problem}

\noindent Sarrus a remarqué la traîtrise de 341 en 1819. Les trompeurs universels de
la question (c) sont les nombres de Carmichael, caractérisés par Korselt en 1899, onze ans
avant que Carmichael ne publie le plus petit d'entre eux, 561, en 1910, et en 1994 Alford, Granville et Pomerance ont démontré
qu'il en existe une infinité : le test de Fermat n'est donc pas seulement percé, il l'est
irrémédiablement. Voilà pourquoi le test de primalité pratique a eu besoin du raffinement
plus puissant de Miller--Rabin, et pourquoi l'existence d'un test rapide et \textit{certain}
(AKS, 2002 ; voyez PRF-20) fut un véritable événement.

\begin{refs}
\item Contexte : Nombre pseudo-premier --- Wikipédia (\url{https://fr.wikipedia.org/wiki/Nombre_pseudo-premier})
\item Contexte : Nombre de Carmichael --- Wikipédia (\url{https://fr.wikipedia.org/wiki/Nombre_de_Carmichael})
\item Article : W. R. Alford, A. Granville \& C. Pomerance, \og{} There are infinitely many Carmichael numbers \fg{} (1994), \textit{Annals of Mathematics} 139
\end{refs}

\bigskip

\begin{problem}[{Le théorème de Wilson : un test parfait et inutile --- Licence}]
\textbf{PRF-20.} \textbf{Problème.} Une caractérisation exacte de la primalité --- et l'écart entre un
critère et un algorithme.
\begin{enumerate}
\item[(a)] Démontrez le théorème de Wilson : un entier
$ n >1 $ est premier si et seulement si
\[ (n-1)! \equiv -1 \pmod{n}. \]
Pour le sens \og{} composé \fg{}, démontrez le fait plus précis que $ (n-1)! \equiv 0 \pmod{n} $
pour tout composé $ n >4 $.
\item[(b)] Contrairement au test de Fermat, c'est une
caractérisation exacte de la primalité --- et pourtant elle ne vaut rien en pratique. Rendez
cela précis : estimez le nombre de multiplications nécessaires pour calculer directement
$ (n-1)! \bmod n $, exprimez-le en fonction du nombre de chiffres de n, et comparez au
coût d'un seul test de Fermat calculé par élévations au carré répétées. Concluez en
explicitant la distinction entre un \textit{critère} et un \textit{algorithme}. (On ne connaît
aucune méthode calculant $ (n-1)! \bmod n $ en un temps polynomial en le nombre de
chiffres --- si vous en trouvez une, vous aurez trouvé un nouveau test de primalité.)
\end{enumerate}

\end{problem}

\noindent L'énoncé était connu d'Ibn al-Haytham vers l'an 1000. John Wilson,
étudiant d'Edward Waring à Cambridge, l'a redécouvert ; Waring l'a publié en 1770, sans
démonstration, déplorant que les théorèmes de ce genre parussent hors d'atteinte faute d'une
notation pour les nombres premiers --- ce à quoi Gauss, qui en donna une démonstration dans les
\textit{Disquisitiones}, rétorqua que de telles vérités demandent des notions, non des
notations. Lagrange en a donné la première démonstration en 1771. La morale de la question
(b) est sérieuse : l'écart entre caractériser un objet et le calculer efficacement est
exactement celui que le théorème d'Agrawal, Kayal et Saxena de 2002 (\og{} PRIMES is in P \fg{}) a
fini par combler pour la primalité.

\begin{refs}
\item Contexte : Théorème de Wilson --- Wikipédia (\url{https://fr.wikipedia.org/wiki/Th%C3%A9or%C3%A8me_de_Wilson})
\item Contexte : Test de primalité AKS --- Wikipédia (\url{https://fr.wikipedia.org/wiki/Test_de_primalit%C3%A9_AKS})
\item Cours : MIT OCW 18.781 Theory of Numbers (en anglais) (\url{https://ocw.mit.edu/courses/18-781-theory-of-numbers-spring-2012/})
\end{refs}

\bigskip

\begin{problem}[{L'argument diagonal de Cantor --- Licence}]
\textbf{PRF-21.} \textbf{Problème.} La méthode diagonale de Cantor, parcourue du concret à
l'abstrait.
\begin{enumerate}
\item[(a)] Démontrez qu'il n'existe pas de surjection de $\mathbb{N}$ sur
l'ensemble de toutes les suites infinies de 0 et de 1 : étant donné une liste quelconque
$ s_1, s_2, s_3, \dots $ de telles suites, construisez une suite qui ne figure nulle part
dans la liste. Déduisez-en que les nombres réels ne sont pas dénombrables.
\item[(b)] Adaptez l'argument pour démontrer le théorème de Cantor : aucun
ensemble X n'admet de surjection sur l'ensemble de ses parties $ \mathcal{P}(X) $.
\item[(c)] La construction
de (a) est un canevas. Énoncez-la assez abstraitement pour que le même squelette donne (b),
et nommez l'autoréférence qui en est le cœur.
\end{enumerate}

\end{problem}

\noindent Cantor a d'abord démontré la non-dénombrabilité des réels en 1874 par un
argument différent, moins souple ; la méthode diagonale est apparue en 1891 et a changé la
discipline. Elle fut attaquée en son temps --- Kronecker tenait les infinis de Cantor pour
illégitimes --- et défendue par la fameuse déclaration de Hilbert : \og{} nul ne doit nous chasser
du paradis que Cantor a créé \fg{}. Le squelette diagonal est devenu l'un des canevas de
démonstration les plus lourds de conséquences jamais conçus : le théorème d'incomplétude de
Gödel et la démonstration par Turing de l'indécidabilité du problème de l'arrêt sont l'un et
l'autre des arguments diagonaux. Apprendre à voir le canevas unique sous les trois théorèmes,
voilà le véritable exercice ici.

\begin{refs}
\item Contexte : Argument de la diagonale de Cantor --- Wikipédia (\url{https://fr.wikipedia.org/wiki/Argument_de_la_diagonale_de_Cantor})
\item Lecture : R. Hammack, \textit{Book of Proof}, ch. 14
\end{refs}

\bigskip

\begin{problem}[{e est irrationnel --- Licence}]
\textbf{PRF-22.} \textbf{Problème.} Soit $ e = \sum_{n=0}^{\infty} \frac{1}{n!} $.
\begin{enumerate}
\item[(a)] À l'aide de la série, démontrez que e est irrationnel. (Le procédé classique : en
supposant $ e = p/q $, examinez $ q! \, e $ et piégez une certaine quantité strictement
entre 0 et 1.)
\item[(b)] Démontrez que $ e^2 $ est irrationnel.
\item[(c)] Expliquez pourquoi votre argument de (a) démontre quelque chose d'un peu plus fort que
l'irrationalité --- à quelle vitesse des rationnels peuvent-ils approcher e ? --- et pourquoi il
ne dit néanmoins rien sur, disons, $ e + \pi $.
\end{enumerate}

\end{problem}

\noindent Euler a démontré l'irrationalité de e en 1737 par sa fraction continue ; le
court argument par les séries demandé ici est d'ordinaire attribué à Fourier (1815).
Liouville a poussé jusqu'à $ e^2 $ et au-delà en 1840, et Hermite a démontré la
transcendance de e en 1873 --- le premier nombre apparu naturellement dont on l'ait démontré.
La question (c) est une leçon sur les limites d'une technique : à ce jour personne ne sait
démontrer que $ e + \pi $ est irrationnel, rappel qu'en mathématiques la distance entre
\og{} manifestement vrai \fg{} et \og{} démontré \fg{} peut être un abîme. Le chapitre \og{} Some irrational
numbers \fg{} d'Aigner et Ziegler mène cette histoire bien plus loin.

\begin{refs}
\item Contexte : Démonstration de l'irrationalité de e --- Wikipédia (en anglais) (\url{https://en.wikipedia.org/wiki/Proof_that_e_is_irrational})
\item Lecture : M. Aigner \& G. M. Ziegler, \textit{Proofs from THE BOOK}, ch. \og{} Some irrational numbers \fg{}
\end{refs}

\bigskip

\begin{problem}[{L'inégalité arithmético-géométrique par récurrence avant-arrière --- Licence}]
\textbf{PRF-23.} \textbf{Problème.} Démontrez l'inégalité entre moyenne arithmétique et moyenne
géométrique : pour des réels positifs ou nuls $ a_1, \dots, a_n $,
\[ \sqrt[n]{a_1 a_2 \cdots a_n} \;\le\; \frac{a_1 + a_2 + \cdots + a_n}{n}, \]
avec égalité exactement lorsque tous les $ a_i $ sont égaux --- par la récurrence
avant-arrière de Cauchy :
\begin{enumerate}
\item[(a)] démontrez le cas $ n = 2 $ à la main ;
\item[(b)] montrez que le cas n entraîne le cas
2n (bond en avant à travers les puissances de deux) ;
\item[(c)] montrez que le cas n entraîne le cas
$ n - 1 $ (pas en arrière, par un choix astucieux de la variable écartée) ;
\item[(d)] expliquez
soigneusement pourquoi (a)--(c) couvrent tous les n, et suivez à la trace où loge le cas
d'égalité à chaque étape. Pourquoi la récurrence ordinaire peine-t-elle ici, tandis que
cette étrange démarche réussit ?
\end{enumerate}

\end{problem}

\noindent C'est la démonstration propre à Cauchy, tirée du \textit{Cours d'analyse}
(1821), et le schéma de récurrence est aussi célèbre que l'inégalité : il grimpe vers
l'infini en doublant, puis redescend un barreau à la fois. L'inégalité
arithmético-géométrique a accumulé des dizaines de démonstrations --- Pólya disait que sa
préférée (par l'inégalité $ e^x \ge 1 + x $) lui était venue en rêve --- et la
\textit{Cauchy-Schwarz Master Class} de Steele en fait la deuxième halte d'une visite de tout
le paysage des inégalités.

\begin{refs}
\item Contexte : Inégalité arithmético-géométrique --- Wikipédia (\url{https://fr.wikipedia.org/wiki/In%C3%A9galit%C3%A9_arithm%C3%A9tico-g%C3%A9om%C3%A9trique})
\item Lecture : J. M. Steele, \textit{The Cauchy-Schwarz Master Class} (MAA/Cambridge, 2004), ch. 2
\end{refs}

\bigskip

\begin{problem}[{Trois démonstrations de Cauchy--Schwarz --- Licence}]
\textbf{PRF-24.} \textbf{Problème.} Démontrez que, pour tous réels
$ a_1, \dots, a_n, b_1, \dots, b_n $,
\[ \left( \sum_{i=1}^{n} a_i b_i \right)^{\!2} \;\le\; \left( \sum_{i=1}^{n} a_i^2 \right) \left( \sum_{i=1}^{n} b_i^2 \right), \]
par trois arguments véritablement différents. Familles candidates : un argument de polynôme
du second degré et de discriminant ; une identité algébrique exacte exhibant la différence
des deux membres comme une somme de carrés ; une récurrence sur n ; une normalisation qui
ramène le tout à l'inégalité arithmético-géométrique.
\begin{enumerate}
\item[(a)] Rédigez les trois démonstrations, et pour chacune déterminez le cas d'égalité qu'elle
fournit.
\item[(b)] Jugez ensuite : laquelle de vos démonstrations survit, mot pour mot, lorsque les sommes
sont remplacées par des intégrales ?
\item[(c)] Laquelle survit dans un espace préhilbertien abstrait, où les coordonnées n'existent
pas ?
\end{enumerate}

\end{problem}

\noindent Cauchy a démontré l'inégalité pour les sommes en 1821 ; Bouniakovski est
passé aux intégrales en 1859 ; Schwarz a redécouvert la forme intégrale en 1888 en bâtissant
la géométrie des espaces de fonctions, ce qui explique que la version abstraite porte son
nom. C'est peut-être l'inégalité la plus utilisée des mathématiques --- toute la géométrie des
espaces de Hilbert s'y appuie --- et Steele ouvre son livre exactement sur cet exercice, parce
qu'un résultat aussi central mérite d'être compris depuis plusieurs directions à la fois.
Trouver des démonstrations \og{} véritablement différentes \fg{}, et dire ce que cela signifie, c'est
la moitié du problème.

\begin{refs}
\item Contexte : Inégalité de Cauchy-Schwarz --- Wikipédia (\url{https://fr.wikipedia.org/wiki/In%C3%A9galit%C3%A9_de_Cauchy-Schwarz})
\item Lecture : J. M. Steele, \textit{The Cauchy-Schwarz Master Class} (MAA/Cambridge, 2004), ch. 1
\end{refs}

\bigskip

\begin{problem}[{La formule d'Euler pour les polyèdres --- Licence}]
\textbf{PRF-25.} \textbf{Problème.} La formule d'Euler pour les polyèdres, et le contre-exemple qui met à
l'épreuve chacune de ses démonstrations.
\begin{enumerate}
\item[(a)] Démontrez que tout graphe planaire connexe à V
sommets, E arêtes et F faces (en comptant la face extérieure non bornée) vérifie
\[ V - E + F = 2, \]
et déduisez-en l'énoncé classique pour les polyèdres convexes.
\item[(b)] Calculez $ V - E + F $
pour un polyèdre en \og{} cadre de tableau \fg{} --- un solide percé de part en part d'un trou
rectangulaire --- et identifiez exactement quelle hypothèse de votre démonstration il
viole.
\item[(c)] La formule d'Euler a de
nombreuses démonstrations, dont certaines subtilement incomplètes ; écrivez la forme la plus
forte de l'hypothèse de planarité que votre propre démonstration utilise réellement, et
expliquez comment vous la défendriez face à un sceptique armé de la question (b).
\end{enumerate}

\end{problem}

\noindent Euler a annoncé la formule dans une lettre à Goldbach en novembre 1750,
étonné que \og{} ces propriétés générales de la géométrie des solides n'aient pas été remarquées
plus tôt \fg{} ; sa propre démonstration avait des lacunes, Legendre en donna une correcte par la
géométrie sphérique en 1794, et la démonstration de Cauchy de 1811 introduisit le point de vue
des graphes planaires. L'histoire de contre-exemples et de réparations de ce théorème --- les
exceptions de Lhuilier, dont le cadre de tableau --- fait tout le sujet de \textit{Proofs and
Refutations} de Lakatos (1976), le livre le plus célèbre jamais écrit sur la façon dont
les démonstrations évoluent réellement : définition et théorème négociant l'un avec l'autre
jusqu'à ce que tous deux tiennent bon.

\begin{refs}
\item Contexte : Caractéristique d'Euler --- Wikipédia (\url{https://fr.wikipedia.org/wiki/Caract%C3%A9ristique_d%27Euler})
\item Lecture : I. Lakatos, \textit{Proofs and Refutations} (Cambridge, 1976)
\item Lecture : M. Aigner \& G. M. Ziegler, \textit{Proofs from THE BOOK}
\end{refs}

\bigskip

\begin{problem}[{Sylvester--Gallai --- Licence}]
\textbf{PRF-26.} \textbf{Problème.} Un théorème de géométrie réelle, et la configuration complexe qui
montre qu'il parle bel et bien de $\mathbb{R}$.
\begin{enumerate}
\item[(a)] Démontrez le théorème de Sylvester--Gallai : tout ensemble
fini de points du plan, non tous alignés, détermine au moins une \textit{droite
ordinaire} --- une droite passant par exactement deux des points.
\item[(b)] Montrez que
l'énoncé tombe en défaut sur les nombres complexes : il existe un ensemble fini de points
du plan projectif complexe, non tous alignés, dans lequel toute droite passant par deux des
points en contient un troisième. (Une célèbre configuration de neuf points fait
l'affaire.)
\end{enumerate}

\end{problem}

\noindent Sylvester a posé (a) dans l'\textit{Educational Times} en 1893 ; la question
a dormi cinquante ans, jusqu'à ce qu'Erdős la réveille en 1943, et Gallai l'a démontrée en
1944. Melchior avait, sans qu'on le remarque, démontré une version projective duale en 1941.
La démonstration à traquer est celle de L. M. Kelly, un argument extrémal de deux lignes (la
technique de PRF-13) qu'Erdős avait rangé dans son \og{} Livre \fg{} imaginaire des démonstrations
parfaites ; Aigner et Ziegler l'ont dûment réimprimée. La question (b) --- la configuration
hessienne, faite des neuf points d'inflexion d'une courbe cubique --- montre que le théorème
porte véritablement sur la géométrie réelle, et non sur la seule incidence. L'histoire se
poursuit au niveau recherche dans PRF-36.

\begin{refs}
\item Contexte : Théorème de Sylvester-Gallai --- Wikipédia (\url{https://fr.wikipedia.org/wiki/Th%C3%A9or%C3%A8me_de_Sylvester-Gallai})
\item Contexte : Configuration hessienne --- Wikipédia (\url{https://fr.wikipedia.org/wiki/Configuration_hessienne})
\item Lecture : M. Aigner \& G. M. Ziegler, \textit{Proofs from THE BOOK}
\end{refs}

\bigskip

\begin{problem}[{Réécriture : une démonstration correcte, mal racontée --- Licence}]
\textbf{PRF-27.} \textbf{Problème.} La démonstration suivante de l'énoncé (vrai)
\og{} $ \frac{2n+1}{n+3} \to 2 $ quand $ n \to \infty $ \fg{} est logiquement récupérable mais
mal écrite. \textit{Démonstration : $ \left| \frac{2n+1}{n+3} - 2 \right| = \left| \frac{2n+1-2n-6}{n+3} \right| = \frac{5}{n+3} <\varepsilon $,
donc $ n >\frac{5}{\varepsilon} - 3 $, donc ça marche pour n assez grand, donc on prend
N plus grand que ça et c'est réglé pour tous les n d'après. }~
\begin{enumerate}
\item[(a)] Dressez la liste de tous les défauts que vous pouvez trouver : symboles employés avant
d'être introduits ou quantifiés, raisonnement présenté à l'envers de son ordre logique,
étapes injustifiées, pronoms sans antécédent, conclusions jamais réellement énoncées.
\item[(b)] Réécrivez la démonstration
aux normes d'une revue : la définition de la limite d'abord, N exhibé explicitement, chaque
inégalité s'enchaînant vers l'avant.
\item[(c)] Tirez de l'exercice trois règles de bonne
écriture mathématique, et confrontez-les aux recommandations de Su et à l'essai de
Halmos.
\end{enumerate}

\end{problem}

\noindent La plupart des démonstrations que rencontrent les étudiants sont déjà
polies, si bien que le métier du polissage n'est jamais vu à l'œuvre. L'essai de Halmos de
1970, \og{} How to Write Mathematics \fg{} --- toujours ce que l'on a écrit de meilleur sur le sujet ---
insiste sur le fait qu'une démonstration est un morceau de prose d'exposition dont la tâche
est de convaincre et d'éclairer, non la transcription des brouillons de l'auteur. Transformer
par la réécriture une mauvaise démonstration en une bonne, la logique restant fixée, isole le
savoir-faire de l'écriture de celui de la découverte ; c'est le remède connu le plus rapide à
la soupe de quantificateurs.

\begin{refs}
\item Contexte : Démonstration (logique et mathématiques) --- Wikipédia (\url{https://fr.wikipedia.org/wiki/D%C3%A9monstration_%28logique_et_math%C3%A9matiques%29})
\item Article : P. R. Halmos, \og{} How to Write Mathematics \fg{} (1970), \textit{L'Enseignement Mathématique} 16
\item Article : F. E. Su, \og{} Guidelines for Good Mathematical Writing \fg{} (2015), PDF (en anglais) (\url{https://math.hmc.edu/su/wp-content/uploads/sites/10/2020/08/Guidelines-for-Good-Mathematical-Writing.pdf})
\item Lecture : D. J. Velleman, \textit{How to Prove It}
\end{refs}

\bigskip

\begin{problem}[{Bon ordre et récurrence : un principe, deux habits --- Licence}]
\textbf{PRF-41.} \textbf{Problème.} Le principe du bon ordre (PBO) dit : toute partie non vide d'entiers
strictement positifs possède un plus petit élément. Le principe de récurrence (PR) dit : si
une partie S d'entiers strictement positifs contient 1, et contient $ n+1 $ dès qu'elle
contient n, alors S contient tous les entiers strictement positifs.
\begin{enumerate}
\item[(a)] En admettant le PBO, démontrez le PR.
\item[(b)] En admettant le PR, démontrez le PBO. (Attention au piège classique : la tentative
naturelle a secrètement besoin de la récurrence forte --- soit déduisez d'abord la récurrence
forte du PR, soit trouvez le bon ensemble sur lequel raisonner par récurrence.)
\item[(c)] Montrez que les deux principes tombent en défaut pour les rationnels positifs ou nuls,
et identifiez quel trait structurel de $\mathbb{N}$ chaque défaillance met au jour.
\item[(d)] Dans les axiomes de Peano, la récurrence est l'axiome et le bon ordre est un théorème.
Discutez en un court paragraphe : étant donné l'équivalence, dire de l'un des deux qu'il est
\og{} plus fondamental \fg{}, est-ce une affirmation mathématique ou une affirmation
esthétique ?
\end{enumerate}

\end{problem}

\noindent Cette page spécule sur cette équivalence depuis le début --- PRF-12
échangeait la récurrence forte contre un plus petit contre-exemple, PRF-13 appelait le
principe extrémal le bon ordre en tenue de travail, et la descente de PRF-31 est le bon ordre
parcouru à l'envers. Ce problème acquitte la dette et démontre le taux de change. Le soin
qu'exige la question (b) est toute la leçon : les équivalences entre principes quasi évidents
sont exactement là où se cache le raisonnement circulaire, et la défense de celui qui rédige
est d'afficher, à chaque étape, quel principe il a en main et lequel lui est interdit.
Dedekind (1888) et Peano (1889) ont fixé la récurrence comme axiome ; lequel des deux habits
est le \og{} fondamental \fg{} se discute depuis lors.

\begin{refs}
\item Contexte : Principe du bon ordre --- Wikipédia (en anglais) (\url{https://en.wikipedia.org/wiki/Well-ordering_principle})
\item Contexte : Axiomes de Peano --- Wikipédia (\url{https://fr.wikipedia.org/wiki/Axiomes_de_Peano})
\item Lecture : D. J. Velleman, \textit{How to Prove It}, ch. 6
\item Lecture : R. Hammack, \textit{Book of Proof}, ch. 10
\end{refs}

\bigskip

\begin{problem}[{Le criminel minimal : six couleurs suffisent --- Licence}]
\textbf{PRF-42.} \textbf{Problème.} Une démonstration par contre-exemple minimal s'ouvre ainsi :
\og{} supposons que le théorème soit en défaut, et parmi tous les contre-exemples fixons-en un
de taille minimale \fg{}. Le bénéfice est que toute instance plus petite obéit au théorème et
peut être employée librement ; le but est de fabriquer un contre-exemple plus petit, ce qui
contredit la minimalité.
\begin{enumerate}
\item[(a)] À l'aide de la formule d'Euler (PRF-25), démontrez que tout graphe planaire simple
possède un sommet de degré au plus 5.
\item[(b)] Démontrez par contre-exemple minimal que tout graphe planaire simple admet une
coloration propre en 6 couleurs (des sommets adjacents recevant toujours des couleurs
différentes).
\item[(c)] Réécrivez la démonstration de (b) comme une récurrence sur le nombre de sommets, et
énoncez précisément comment les deux tournures se traduisent l'une dans l'autre (comparez
PRF-12 et PRF-41).
\item[(d)] À partir des références, cartographiez ce qui se produit quand le compte diminue :
l'argument des chaînes de Kempe donne 5 couleurs honnêtement --- étudiez-le --- et a donné 4
couleurs de façon fallacieuse pendant onze ans. Où exactement la démonstration de Kempe
cassait-elle ?
\end{enumerate}

\end{problem}

\noindent \og{} Criminel minimal \fg{} est le nom affectueux de la technique, et la
coloration de graphes en est le foyer historique. Kempe a publié une démonstration du
théorème des quatre couleurs en 1879 ; Heawood en a trouvé la faille en 1890 et a sauvé du
naufrage le théorème des cinq couleurs --- la démonstration ratée la plus instructive des
mathématiques, une suite grandeur nature à PRF-05 et PRF-06. Le cas des quatre couleurs a
finalement été réglé par Appel et Haken en 1976, en le ramenant à une énorme vérification de
cas menée par ordinateur, ce qui a allumé sur ce qu'est une démonstration un débat dont
l'essai de Thurston (PRF-32) porte encore l'écho ; la formalisation vérifiée par machine de
Gonthier en Coq (2005) a bouclé la boucle par l'autre bout.

\begin{refs}
\item Contexte : Théorème des cinq couleurs --- Wikipédia (\url{https://fr.wikipedia.org/wiki/Th%C3%A9or%C3%A8me_des_cinq_couleurs})
\item Contexte : Théorème des quatre couleurs --- Wikipédia (\url{https://fr.wikipedia.org/wiki/Th%C3%A9or%C3%A8me_des_quatre_couleurs})
\item Lecture : M. Aigner \& G. M. Ziegler, \textit{Proofs from THE BOOK}, ch. \og{} Five-coloring plane graphs \fg{}
\item Article : G. Gonthier, \og{} Formal Proof --- The Four-Color Theorem \fg{} (2008), \textit{Notices of the AMS} 55
\end{refs}

\bigskip

\begin{problem}[{L'atelier du contre-exemple --- Licence}]
\textbf{PRF-43.} \textbf{Problème.} Pour tuer un énoncé universel, un seul objet bien fait suffit. Pour
chacune des affirmations ci-dessous, démontrez-la ou réfutez-la ; quand vous réfutez,
présentez le contre-exemple avec le soin dû à un théorème --- l'objet, la vérification, la
conclusion.
\begin{enumerate}
\item[(a)] Si $ f\colon \mathbb{R} \to \mathbb{R} $ et si $ |f| $ est continue, alors f est
continue.
\item[(b)] Si $ \sum a_n $ converge, alors $ \sum a_n^2 $ converge.
\item[(c)] Si $ f\colon \mathbb{R} \to \mathbb{R} $ est dérivable partout, alors sa dérivée
$ f' $ est continue.
\item[(d)] Toute fonction $ f\colon \mathbb{R} \to \mathbb{R} $ vérifiant
$ f(x+y) = f(x) + f(y) $ pour tous x, y est de la forme $ f(x) = cx $. (Celle-ci est
particulière : tranchez-la, puis discutez de ce que \og{} construire un contre-exemple \fg{} peut
vouloir dire quand --- comme ici, par un théorème de Hamel --- un contre-exemple existe sans
qu'aucun puisse être écrit explicitement, mais seulement invoqué à partir de l'axiome du
choix.)
\end{enumerate}

\end{problem}

\noindent Les contre-exemples ne sont pas le contraire des démonstrations, mais une
espèce de démonstrations, et en construire un est un problème de conception : dressez la
liste des exigences, trouvez leurs tensions, et bâtissez l'objet exigence par exigence.
L'analyse en est la carrière classique --- \textit{Counterexamples in Analysis} de Gelbaum et
Olmsted en est tout un musée --- parce que ses définitions sont juste assez permissives pour
admettre des monstres. La question (d) ouvre une porte plus profonde : Hamel a montré en 1905
que des fonctions additives sauvages existent, et pourtant chacune d'elles est non mesurable,
de sorte qu'aucune formule n'en exhibera jamais une. Ce qu'\textit{est} un contre-exemple,
lorsqu'il n'existe que par l'axiome du choix, se poursuit en PRF-45 et PRF-48.

\begin{refs}
\item Contexte : Contre-exemple --- Wikipédia (\url{https://fr.wikipedia.org/wiki/Contre-exemple})
\item Contexte : Équation fonctionnelle de Cauchy --- Wikipédia (\url{https://fr.wikipedia.org/wiki/%C3%89quation_fonctionnelle_de_Cauchy})
\item Lecture : B. R. Gelbaum \& J. M. H. Olmsted, \textit{Counterexamples in Analysis} (1964)
\item Lecture : I. Lakatos, \textit{Proofs and Refutations} (Cambridge, 1976)
\end{refs}

\bigskip

\begin{problem}[{La négation sans \og{} non \fg{} --- Licence, short}]
\textbf{PRF-44.} \textbf{Problème.} Écrivez la négation exacte de \og{} pour tout
$ \varepsilon >0 $ il existe un $ \delta >0 $ tel que
$ |x - a| <\delta $ implique $ |f(x) - f(a)| <\varepsilon $ \fg{} --- quantificateurs
dans le bon ordre, sans le mot \og{} non \fg{} et sans aucun symbole de négation. Puis maniez-la :
démontrez que la fonction valant $ f(x) = 1 $ pour x rationnel et $ f(x) = 0 $ pour x
irrationnel n'est continue en aucun point.
\end{problem}

\noindent Faire traverser à une négation une file de quantificateurs --- chaque $\forall$
basculant en $\exists$ et retour, l'implication finale s'effondrant en un témoin --- est la
gymnastique de l'analyse ; la plupart des démonstrations de limite ratées sont des erreurs
de quantificateurs déguisées. La fonction test est celle de Dirichlet (1829), la première
fonction célèbre discontinue partout, et démontrer sa discontinuité est exactement un jeu
consistant à produire les bons témoins.

\begin{refs}
\item Contexte : Fonction continue nulle part --- Wikipédia (\url{https://fr.wikipedia.org/wiki/Fonction_continue_nulle_part})
\item Lecture : D. J. Velleman, \textit{How to Prove It}, ch. 2
\end{refs}

\bigskip

\begin{problem}[{Un irrationnel élevé à une puissance irrationnelle --- Licence, short}]
\textbf{PRF-45.} \textbf{Problème.} Démontrez qu'il existe des nombres irrationnels a et
b tels que $ a^b $ soit rationnel. (Un célèbre argument à deux cas y parvient sans jamais
déterminer lequel de ses deux candidats convient ; trouvez-le --- ou toute autre démonstration
--- puis dites clairement si votre démonstration exhibe un couple effectif.)
\end{problem}

\noindent Le spécimen classique de la démonstration d'existence non constructive :
le tour à deux cas, appliqué à $ \sqrt{2}^{\sqrt{2}} $, convainc pleinement tout en vous
laissant incapable de nommer le couple qu'il promet --- tout repose sur le principe du tiers
exclu, que les mathématiciens constructivistes refusent d'accorder. Le théorème de
Gelfond--Schneider (1934) règle l'affaire hors du prétoire, et à grands frais :
$ \sqrt{2}^{\sqrt{2}} $ est en réalité transcendant. Une pièce jumelle des contre-exemples
innommables de PRF-43 et des démonstrations d'existence probabilistes de PRF-48.

\begin{refs}
\item Contexte : Démonstration constructive --- Wikipédia (\url{https://fr.wikipedia.org/wiki/D%C3%A9monstration_constructive})
\item Contexte : Théorème de Gelfond-Schneider --- Wikipédia (\url{https://fr.wikipedia.org/wiki/Th%C3%A9or%C3%A8me_de_Gelfond-Schneider})
\end{refs}

\bigskip

\begin{problem}[{La démonstration d'une seule phrase de Zagier --- Licence}]
\textbf{PRF-47.} \textbf{Problème.} Soit p un nombre premier tel que $ p \equiv 1 \pmod 4 $ et soit
\[ S \;=\; \{\, (x,y,z) \in \mathbb{Z}_{>0}^{3} \;:\; x^{2} + 4yz = p \,\}. \]
Une \textit{involution} d'un ensemble est une application égale à sa propre réciproque ; sur un
ensemble fini, le nombre de ses points fixes a la même parité que le cardinal de
l'ensemble.
\begin{enumerate}
\item[(a)] Montrez que S est fini et non vide, que $ (x,y,z) \mapsto (x,z,y) $ est une
involution de S, et que cette involution possède un point fixe si et seulement si p est une
somme de deux carrés.
\item[(b)] Vérifiez que
\[ (x,y,z) \;\longmapsto\;
\begin{cases}
(x+2z,\; z,\; y-x-z) & \text{si } x <y-z,\\
(2y-x,\; y,\; x-y+z) & \text{si } y-z <x <2y,\\
(x-2y,\; x-y+z,\; y) & \text{si } x >2y
\end{cases} \]
est bien définie sur S --- en particulier que les trois cas sont exhaustifs et mutuellement
exclusifs, et qu'aucune coordonnée ne quitte jamais $ \mathbb{Z}_{>0} $ ---, que c'est
une involution, et qu'elle possède exactement un point fixe.
\item[(c)] Assemblez (a) et (b) en une démonstration du théorème de Fermat : tout nombre premier
$ p \equiv 1 \pmod 4 $ est une somme de deux carrés. Répondez ensuite à la question
inconfortable que soulève cette démonstration : ayant vérifié chaque ligne, qu'avez-vous
appris ? Rédigez un paragraphe sur les raisons pour lesquelles cette démonstration convainc
sans expliquer, et comparez-la à une démonstration passant par les entiers de Gauss
$ \mathbb{Z}[i] $ ou par le lemme de Thue.
\end{enumerate}

\end{problem}

\noindent Fermat a annoncé le théorème à Mersenne dans une lettre de décembre 1640
et, comme à son habitude, a gardé la démonstration pour lui ; il a fallu à Euler sept ans
d'efforts intermittents avant d'en publier une en 1749. En 1990 Don Zagier a comprimé un
argument de Heath-Brown en une seule phrase imprimée dans l'\textit{American Mathematical
Monthly} --- très littéralement une phrase, l'application affichée ci-dessus plus une
proposition subordonnée --- et c'est depuis lors la pièce à conviction standard dans la
dispute sur ce à quoi \textit{sert} une démonstration. Chaque étape est vérifiable en un
après-midi par un étudiant appliqué ; d'où vient l'application, cela reste invisible. Des
exposés ultérieurs restituent une part du sens en dessinant les triplets comme des figures
en \og{} moulin à vent \fg{}, un carré muni de quatre bras rectangulaires, moyennant quoi les trois
cas deviennent trois façons de rebâtir le moulin. Comparez l'involution d'ici à la stratégie
du miroir de PRF-37 : toutes deux laissent une symétrie faire le dénombrement.

\begin{refs}
\item Contexte : Théorème des deux carrés de Fermat --- Wikipédia (\url{https://fr.wikipedia.org/wiki/Th%C3%A9or%C3%A8me_des_deux_carr%C3%A9s_de_Fermat})
\item Contexte : Involution (mathématiques) --- Wikipédia (\url{https://fr.wikipedia.org/wiki/Involution_%28math%C3%A9matiques%29})
\item Article : D. Zagier, \og{} A One-Sentence Proof That Every Prime \textit{p} $\equiv$ 1 (mod 4) Is a Sum of Two Squares \fg{} (1990), \textit{American Mathematical Monthly} 97, p. 144
\item Lecture : M. Aigner \& G. M. Ziegler, \textit{Proofs from THE BOOK}, le chapitre sur la représentation des nombres comme sommes de deux carrés
\end{refs}

\bigskip

\begin{problem}[{Absurde ou construction : ce qu'une démonstration vous remet --- Licence}]
\textbf{PRF-56.} \textbf{Problème.} Un raisonnement par l'absurde suppose P fausse, en tire une absurdité,
et conclut P. Une démonstration directe de P ne mentionne jamais la négation de P. Les deux
ne sont pas la même chose, et la différence mérite d'être suivie à la trace.
\begin{enumerate}
\item[(a)] Chacun des énoncés suivants est habituellement présenté comme un raisonnement par
l'absurde. Pour chacun, décidez si l'absurde fait véritablement le travail ou si l'argument
est une démonstration directe (ou une contraposée) déguisée, et réécrivez-le sous la forme
correcte la plus dépouillée : (i) $ \sqrt{2} $ est irrationnel (PRF-02) ; (ii) si
$ n^2 $ est pair, alors n est pair ; (iii) il existe une infinité de nombres premiers
(PRF-14) ; (iv) une fonction strictement croissante
$ f\colon \mathbb{R} \to \mathbb{R} $ est injective.
\item[(b)] Explicitez la logique. Montrez que la démonstration d'un énoncé \textit{négatif} \og{} non Q \fg{}
obtenue en supposant Q et en aboutissant à une absurdité est une démonstration directe de
\og{} non Q \fg{} et n'emploie aucun principe supplémentaire ; montrez ensuite que la démonstration
d'un énoncé \textit{positif} P obtenue en supposant \og{} non P \fg{} et en aboutissant à une
absurdité emploie exactement l'inférence de \og{} non non P \fg{} à P. Identifiez à quelle loi
classique cette inférence est équivalente.
\item[(c)] Voici maintenant le prix. Une démonstration directe de \og{} il existe x tel que
$ Q(x) $ \fg{} vous remet d'ordinaire un x ; un raisonnement par l'absurde n'y est pas tenu.
Prenez l'argument à deux cas de PRF-45 et la coloration aléatoire de PRF-48, et dans chaque
cas désignez l'unique étape qui détruit le témoin.
\item[(d)] Enfin, lisez ce que les constructivistes font à la place. Énoncez la lecture de
Brouwer--Heyting--Kolmogorov du \og{} ou \fg{} et du \og{} il existe \fg{}, expliquez pourquoi, sous cette
lecture, l'argument de PRF-45 n'est pas du tout une démonstration, puis dites précisément ce
que l'on y gagne : nommez une chose qu'une démonstration constructive d'un énoncé
d'existence vous donne et qu'une démonstration classique ne donne pas.
\end{enumerate}

\end{problem}

\noindent Brouwer a commencé à attaquer le principe du tiers exclu sans restriction
vers 1908, et Hilbert a répondu qu'ôter ce principe au mathématicien reviendrait à refuser le
télescope à l'astronome ou à interdire au boxeur l'usage de ses poings. La querelle n'a jamais
été tranchée, plutôt affinée : Gödel et Gentzen ont montré en 1933 que l'arithmétique
classique se traduit dans l'arithmétique intuitionniste par insertion de doubles négations,
de sorte que le raisonnement classique ne démontre rien de neuf quant à la cohérence --- ce
qu'il change, c'est le \textit{contenu} d'une démonstration, non sa sûreté. Les \textit{Foundations
of Constructive Analysis} d'Errett Bishop (1967) ont ensuite fait ce que les arguments ne
peuvent pas faire : reconstruire constructivement une large part de l'analyse, pour montrer
que c'était faisable. L'enjeu n'est plus seulement philosophique. Sous la correspondance de
Curry--Howard, une démonstration constructive d'un énoncé d'existence \textit{est} un programme
qui produit l'objet, ce qui explique que les assistants de preuve se soucient intensément de
savoir quelles étapes sont classiques.

\begin{refs}
\item Contexte : Principe du tiers exclu --- Wikipédia (\url{https://fr.wikipedia.org/wiki/Principe_du_tiers_exclu})
\item Contexte : Logique intuitionniste --- Wikipédia (\url{https://fr.wikipedia.org/wiki/Logique_intuitionniste})
\item Lecture : E. Bishop, \textit{Foundations of Constructive Analysis} (McGraw-Hill, 1967)
\item Lecture : D. Bridges \& F. Richman, \textit{Varieties of Constructive Mathematics} (Cambridge, 1987)
\end{refs}

\bigskip

\begin{problem}[{Récurrence sur la structure, non sur les nombres --- Licence, short}]
\textbf{PRF-57.} \textbf{Problème.} Définissez un arbre binaire comme étant soit une
feuille unique, soit un couple ordonné d'arbres binaires, et définissez une formule
propositionnelle comme étant soit un atome, soit $ \neg \varphi $ pour une formule $\varphi$, soit
$ (\varphi \wedge \psi) $ pour des formules $\varphi$ et $\psi$. Démontrez par induction structurelle
que tout arbre binaire a exactement une feuille de plus qu'il n'a de nœuds internes, et que
toute formule propositionnelle contient autant de parenthèses ouvrantes que de fermantes ---
puis énoncez sur quelle relation bien fondée chacune de ces inductions s'appuie réellement,
et pourquoi la récurrence ordinaire sur les entiers naturels en est un cas particulier
plutôt qu'une rivale.
\end{problem}

\noindent On enseigne d'ordinaire la récurrence comme un énoncé sur $\mathbb{N}$, ce qui masque
ce qui la fait fonctionner : n'importe quelle relation sans chaîne descendante infinie fera
l'affaire, et les objets construits par une grammaire finie portent gratuitement une telle
relation. R. M. Burstall a mis ce point au service de l'informatique en 1969, et l'induction
structurelle est aujourd'hui la manière dont on vérifie compilateurs, systèmes de types et
assistants de preuve --- les démonstrations de mathlib sur la syntaxe sont des inductions
structurelles presque à la ligne près. Le fait jumeau, à savoir que les \textit{définitions}
récursives sur de telles structures sont légitimes, est un théorème à part entière et vaut
d'être poursuivi après celui-ci. Comparez avec PRF-41, qui acquitte la dette entre récurrence
et bon ordre sur $\mathbb{N}$.

\begin{refs}
\item Contexte : Induction structurelle --- Wikipédia (\url{https://fr.wikipedia.org/wiki/Induction_structurelle})
\item Contexte : Relation bien fondée --- Wikipédia (\url{https://fr.wikipedia.org/wiki/Relation_bien_fond%C3%A9e})
\item Cours : MIT OCW 6.042J Mathematics for Computer Science (en anglais) (\url{https://ocw.mit.edu/courses/6-042j-mathematics-for-computer-science-fall-2010/})
\end{refs}

\bigskip

\begin{problem}[{La démonstration par la figure, rendue rigoureuse : l'équidécomposabilité --- Licence}]
\textbf{PRF-58.} \textbf{Problème.} Dites que deux polygones sont \textit{équidécomposables} si l'un peut
être découpé en un nombre fini de pièces polygonales que l'on peut réassembler, par
translations et rotations, en l'autre. PRF-11 demandait quand une figure de découpage compte
comme une démonstration ; ce problème répond à la question en dimension deux, puis montre où
les figures s'épuisent --- et le démontre.
\begin{enumerate}
\item[(a)] Démontrez le théorème de Wallace--Bolyai--Gerwien : deux polygones sont
équidécomposables si et seulement s'ils ont la même aire. (Une voie praticable : tout
polygone se découpe en triangles, tout triangle en un rectangle, tout rectangle en un
rectangle de largeur prescrite quelconque, et l'équidécomposabilité est une relation
d'équivalence.)
\item[(b)] Auditez ce que vous avez fait à l'aune de PRF-11. Lesquelles des hypothèses qu'une
figure de découpage formule en silence votre démonstration a-t-elle désormais acquittées, et
lesquelles sont-elles encore supposées --- que les pièces ne se rencontrent que le long de
leurs bords, que l'aire est finiment additive, que les déplacements préservent l'aire ?
Énoncez comme hypothèses celles que vous conservez.
\item[(c)] Passons à la dimension trois. Énoncez le troisième problème de Hilbert, définissez
l'invariant de Dehn d'un polyèdre, démontrez qu'il est additif par découpage, et calculez-le
pour un cube et pour un tétraèdre régulier --- tranchant ainsi la question de savoir si un
cube et un tétraèdre régulier de même volume peuvent être équidécomposables.
\item[(d)] Énoncez le théorème de Sydler (1965) sur les couples de polyèdres qui \textit{sont}
équidécomposables, et dites ce qu'il signifie pour le programme des démonstrations par la
figure dans l'espace. Cherchez ensuite le paradoxe de Banach--Tarski et expliquez, en un
paragraphe, pourquoi il n'est un contre-exemple à rien de ce qui précède.
\end{enumerate}

\end{problem}

\noindent Wallace a démontré le théorème bidimensionnel en 1807, et Bolyai et
Gerwien l'ont redécouvert indépendamment en 1833 et 1835 ; il dit que dans le plan une
démonstration par découpage n'est jamais l'obstacle --- si deux polygones ont la même aire, la
figure existe et il suffit de la trouver. Gauss s'était plaint dans ses lettres que la théorie
élémentaire du volume parût requérir des arguments d'exhaustion là où l'aire n'en demandait
pas, et Hilbert a converti la plainte en le troisième de ses problèmes de 1900 : exhiber deux
polyèdres de même volume que l'on ne puisse découper l'un en l'autre. Max Dehn a répondu dans
l'année --- le premier des vingt-trois à tomber --- en inventant un invariant bâti sur les
longueurs d'arêtes et les angles dièdres, qui survit au découpage ; Sydler a démontré en 1965
que le volume et l'invariant de Dehn décident ensemble la question complètement, et Jessen a
étendu cela à la dimension quatre. La question analogue dans l'espace hyperbolique et dans
l'espace sphérique est toujours ouverte.

\begin{refs}
\item Contexte : Théorème de Wallace-Bolyai-Gerwien --- Wikipédia (\url{https://fr.wikipedia.org/wiki/Th%C3%A9or%C3%A8me_de_Wallace-Bolyai-Gerwien})
\item Contexte : Troisième problème de Hilbert --- Wikipédia (\url{https://fr.wikipedia.org/wiki/Troisi%C3%A8me_probl%C3%A8me_de_Hilbert})
\item Contexte : Invariant de Dehn --- Wikipédia (en anglais) (\url{https://en.wikipedia.org/wiki/Dehn_invariant})
\item Lecture : M. Aigner \& G. M. Ziegler, \textit{Proofs from THE BOOK}, le chapitre sur le troisième problème de Hilbert
\end{refs}

\bigskip


\section*{Master}

\begin{problem}[{La seconde démonstration d'Euler : la somme des 1/p diverge --- Master}]
\textbf{PRF-28.} \textbf{Problème.} Démontrez que
\[ \sum_{p \text{ premier}} \frac{1}{p} \;=\; \frac{1}{2} + \frac{1}{3} + \frac{1}{5} + \frac{1}{7} + \frac{1}{11} + \cdots \]
diverge.
\begin{enumerate}
\item[(a)] Route 1 : rendez rigoureuse la manipulation d'Euler de 1737 --- reliez
$ \sum_{n \le N} \frac{1}{n} $ au produit
$ \prod_{p \le N} \left(1 - \frac{1}{p}\right)^{-1} $
sur les nombres premiers, prenez les logarithmes, et contrôlez les termes d'erreur.
\item[(b)] Route 2 : donnez la démonstration élémentaire par dénombrement d'Erdős (1938), qui n'a
besoin d'aucun logarithme.
\item[(c)] Tirez ensuite la conséquence qui donne sa force à ce théorème : les nombres premiers
sont en infinité (Euclide, encore !) et, de surcroît, plus \og{} denses \fg{} que les carrés, dont
la somme des inverses converge (PRF-15). Quantifiez cette dernière remarque aussi finement
que vous le pouvez.
\end{enumerate}

\end{problem}

\noindent Les \textit{Variae observationes circa series infinitas} d'Euler (1737) sont
l'acte de naissance de la théorie analytique des nombres : la formule du produit reliant la
fonction zêta aux nombres premiers transforme une question d'arithmétique en une question
d'analyse. Là où la démonstration d'Euclide produit un nouveau nombre premier à la fois,
celle d'Euler mesure combien il y en a --- Mertens a rendu la mesure précise en 1874, montrant
que les sommes partielles croissent comme log log x, peut-être la reptation la plus célèbre
des mathématiques. Un siècle après Euler, Dirichlet a fait tourner le même moteur dans les
progressions arithmétiques, et Riemann a fait passer la formule du produit dans le complexe.
La démonstration alternative d'Erdős est la première entrée de \textit{Proofs from THE
BOOK}.

\begin{refs}
\item Contexte : Série des inverses des nombres premiers --- Wikipédia (\url{https://fr.wikipedia.org/wiki/S%C3%A9rie_des_inverses_des_nombres_premiers})
\item Lecture : M. Aigner \& G. M. Ziegler, \textit{Proofs from THE BOOK}, ch. 1
\item Lecture : G. H. Hardy \& E. M. Wright, \textit{An Introduction to the Theory of Numbers}, ch. 22
\end{refs}

\bigskip

\begin{problem}[{La démonstration topologique de Fürstenberg --- Master}]
\textbf{PRF-29.} \textbf{Problème.} Sur les entiers $\mathbb{Z}$, dites qu'un ensemble est ouvert s'il est une
réunion (éventuellement vide) de progressions arithmétiques bilatères
$ a + b\mathbb{Z} = \{ a + bn : n \in \mathbb{Z} \} $ avec $ b \ne 0 $.
\begin{enumerate}
\item[(a)] Vérifiez que cela définit une topologie sur $\mathbb{Z}$.
\item[(b)] Montrez que tout ouvert non vide est
infini, et que chaque ensemble de base $ a + b\mathbb{Z} $ est fermé autant qu'ouvert.
\item[(c)] Déduisez-en qu'il existe une infinité de nombres premiers, en contemplant l'ensemble
$ \bigcup_{p \text{ premier}} p\mathbb{Z} $ et son complémentaire.
\item[(d)] Vient maintenant la question de
discussion, à laquelle il faut répondre par une position et un argument : est-ce une
démonstration véritablement nouvelle du théorème d'Euclide, ou bien l'arithmétique d'Euclide
elle-même vêtue de langage topologique ? Quel travail, exactement, la topologie
accomplit-elle ?
\end{enumerate}

\end{problem}

\noindent Harry Fürstenberg a publié cette merveille d'une demi-page dans l'\textit{American
Mathematical Monthly} en 1955, alors qu'il était encore étudiant de premier cycle à
l'université Yeshiva ; il devait ensuite révolutionner la théorie ergodique et partager le
prix Abel 2020. Cette démonstration est une favorite précisément parce qu'elle provoque la
question (d) : les mathématiciens ne s'accordent pas sur le point de savoir si la topologie
fait un travail réel ou fournit un emballage brillant, et décider ce qui fait que deux
démonstrations sont \og{} la même \fg{} se révèle une sérieuse question philosophique sur la
démonstration elle-même --- question qui reparaît au niveau recherche en théorie de la
démonstration.

\begin{refs}
\item Contexte : Démonstration de Furstenberg de l'infinité des nombres premiers --- Wikipédia (\url{https://fr.wikipedia.org/wiki/D%C3%A9monstration_de_Furstenberg_de_l%27infinit%C3%A9_des_nombres_premiers})
\item Article : H. Furstenberg, \og{} On the infinitude of primes \fg{} (1955), \textit{American Mathematical Monthly} 62, p. 353
\item Lecture : M. Aigner \& G. M. Ziegler, \textit{Proofs from THE BOOK}, ch. 1
\end{refs}

\bigskip

\begin{problem}[{Schröder--Bernstein --- Master}]
\textbf{PRF-30.} \textbf{Problème.} Le théorème de Schröder--Bernstein : s'il existe des injections
$ f\colon A \to B $ et $ g\colon B \to A $, alors il existe une bijection entre A et
B.
\begin{enumerate}
\item[(a)] Démontrez-le. Votre démonstration ne doit pas invoquer l'axiome du choix. Deux
stratégies classiques à traquer : partitionner A et B en remontant l'ascendance de chaque
élément, en avant et en arrière à travers f et g, et voir où cette ascendance s'arrête ; ou
appliquer un principe de point fixe pour les applications monotones d'ensembles
(Knaster--Tarski).
\item[(b)] Après avoir démontré le théorème, exhibez un couple concret --- les intervalles
$ [0,1] $ et $ (0,1) $, avec des injections de votre choix --- et décrivez explicitement
la bijection que votre démonstration en construit.
\end{enumerate}

\end{problem}

\noindent Un théorème doté d'une comédie de paternité : Cantor l'a énoncé en 1887,
mais sa voie s'appuyait sur des principes équivalents à l'axiome du choix ; Dedekind avait
dans ses carnets, la même année, une démonstration nette et sans axiome du choix, qu'il n'a
jamais publiée ; Schröder a imprimé en 1896 une démonstration que l'on a plus tard trouvée
fallacieuse ; et la première démonstration correcte publiée est venue en 1897 de Felix
Bernstein --- étudiant de dix-neuf ans du séminaire de Cantor --- parvenant à l'impression par le
traité de Borel de 1898. C'est l'outil fondamental qui fait que la comparaison des cardinaux
se comporte comme un ordre, et les deux démonstrations standard offrent un beau contraste de
styles : l'une dynamique et imagée, l'autre relevant de la théorie des ordres et
abstraite.

\begin{refs}
\item Contexte : Théorème de Cantor-Bernstein --- Wikipédia (\url{https://fr.wikipedia.org/wiki/Th%C3%A9or%C3%A8me_de_Cantor-Bernstein})
\item Lecture : P. R. Halmos, \textit{Naive Set Theory} (1960)
\item Lecture : R. Hammack, \textit{Book of Proof}, ch. 14
\end{refs}

\bigskip

\begin{problem}[{La descente infinie de Fermat --- Master}]
\textbf{PRF-31.} \textbf{Problème.} L'arme propre de Fermat : la \textit{descente infinie}.
\begin{enumerate}
\item[(a)] Démontrez par descente infinie que l'équation
$ x^4 + y^4 = z^2 $ n'a pas de solution en entiers strictement positifs : de toute
solution supposée, fabriquez-en une strictement plus petite, et expliquez soigneusement
pourquoi cela constitue une démonstration (quel principe de PRF-12 la garantit ?).
Déduisez-en le cas $ n = 4 $ du grand théorème de Fermat : $ x^4 + y^4 = z^4 $ n'a pas
de solution en entiers strictement positifs.
\item[(b)] Démontrez le théorème de Fermat
sur les triangles rectangles : aucun triangle rectangle dont les côtés ont des longueurs
entières n'a une aire égale à un carré parfait.
\item[(c)] Expliquez le rapport logique entre (a) et (b).
\end{enumerate}

\end{problem}

\noindent Parmi les quarante-huit observations que Fermat a griffonnées dans son
exemplaire de Diophante --- publiées par son fils Samuel en 1670, cinq ans après la mort de
Fermat ---, le théorème sur les triangles rectangles est la seule accompagnée d'une véritable
démonstration : c'est donc l'unique spécimen complet subsistant de la méthode propre du
maître, la \textit{descente infinie ou indéfinie}. La descente, c'est l'absurde mû par le bon
ordre : un plus petit criminel ne peut exister si tout criminel en engendre un plus petit. Le
théorème dit aussi que 1 n'est pas un \textit{nombre congruent} (l'aire d'un triangle rectangle
rationnel), problème que Fibonacci avait soulevé en 1225 ; son prolongement moderne passe par
le théorème de Tunnell (1983) et va droit à la conjecture de Birch et Swinnerton-Dyer ---
voyez la page des problèmes ouverts. L'histoire écrite par Weil reconstitue le raisonnement
de Fermat dans un détail amoureux.

\begin{refs}
\item Contexte : Méthode de descente infinie --- Wikipédia (\url{https://fr.wikipedia.org/wiki/M%C3%A9thode_de_descente_infinie})
\item Contexte : Théorème de Fermat sur les triangles rectangles --- Wikipédia (\url{https://fr.wikipedia.org/wiki/Th%C3%A9or%C3%A8me_de_Fermat_sur_les_triangles_rectangles})
\item Lecture : A. Weil, \textit{Number Theory: An Approach Through History from Hammurapi to Legendre} (1984), ch. II
\end{refs}

\bigskip

\begin{problem}[{Un théorème, trois publics --- Master}]
\textbf{PRF-32.} \textbf{Problème.} Prenez un théorème de cette page que vous avez déjà démontré ---
PRF-02 ou PRF-14 sont idéaux --- et rédigez-en la démonstration trois fois.
\begin{enumerate}
\item[(a)] Pour un lycéen
curieux : aucun symbole au-delà de l'arithmétique, chaque idée motivée, rien de perdu.
\item[(b)] Pour une revue : rigueur complète, dans les conventions de Halmos et de Su, aussi
économe que l'honnêteté le permet.
\item[(c)] Sous forme de squelette de formalisation : dressez la liste de chaque définition,
axiome et lemme qu'un assistant de preuve tel que Lean vous forcerait à fournir --- jusqu'à
la raison pour laquelle un produit d'entiers est un entier --- dans l'ordre des dépendances,
sans nécessairement écrire le code.
\item[(d)] Puis, en un paragraphe : qu'est-ce que chaque public vous a forcé à comprendre que les
deux autres vous laissaient sauter ?
\end{enumerate}

\end{problem}

\noindent L'essai de Thurston \og{} On Proof and Progress in Mathematics \fg{} (1994) ---
écrit en réponse à un débat sur la rigueur dans le \textit{Bulletin of the AMS} --- soutient que
les mathématiques avancent par la compréhension humaine, dont la démonstration formelle est
la partie porteuse mais non l'unique ; une démonstration s'écrit \textit{à} quelqu'un. Le
troisième public n'est plus hypothétique : mathlib, la bibliothèque de la communauté Lean,
contient désormais une large part du canon de premier cycle (théorème d'Euclide compris), et
l'achèvement en 2022 du Liquid Tensor Experiment a montré la démonstration vérifiée par
machine atteignant la frontière de la recherche. Traduire une même vérité dans trois
registres est ce que l'écriture de démonstrations offre de plus proche d'un entraînement
complet.

\begin{refs}
\item Contexte : Lean (assistant de preuve) --- Wikipédia (\url{https://fr.wikipedia.org/wiki/Lean_%28assistant_de_preuve%29})
\item Article : W. P. Thurston, \og{} On Proof and Progress in Mathematics \fg{} (1994), arXiv:math/9404236 (\url{https://arxiv.org/abs/math/9404236})
\item Article : F. E. Su, \og{} Guidelines for Good Mathematical Writing \fg{} (2015), PDF (en anglais) (\url{https://math.hmc.edu/su/wp-content/uploads/sites/10/2020/08/Guidelines-for-Good-Mathematical-Writing.pdf})
\item Lecture : P. R. Halmos, \og{} How to Write Mathematics \fg{} (1970), \textit{L'Enseignement Mathématique} 16
\end{refs}

\bigskip

\begin{problem}[{La méthode probabiliste : démontrer l'existence à pile ou face --- Master}]
\textbf{PRF-48.} \textbf{Problème.} Soit $ R(k,k) $ le plus petit n tel que toute coloration en rouge et
bleu des arêtes du graphe complet $ K_n $ contienne un $ K_k $ monochrome. La méthode
probabiliste démontre qu'un objet existe en montrant qu'un objet tiré au hasard a une
probabilité strictement positive de convenir.
\begin{enumerate}
\item[(a)] Colorez les arêtes de $ K_n $ indépendamment et uniformément au hasard. Démontrez que si
\[ \binom{n}{k} 2^{\,1 - \binom{k}{2}} \;<\; 1 \]
alors $ R(k,k) >n $, et déduisez-en la borne d'Erdős $ R(k,k) >2^{k/2} $ pour
$ k \ge 3 $.
\item[(b)] Le raffinement par suppression : au lieu d'exiger zéro copie monochrome, majorez leur
\textit{nombre espéré}, supprimez un sommet dans chacune, et montrez que cela donne une
minoration strictement meilleure qu'en (a). Énoncez le principe général que vous venez
d'employer, sous une forme qui ne mentionne pas les graphes.
\item[(c)] Énoncez précisément le lemme local de Lovász (Erdős--Lovász, 1975), y compris
l'hypothèse sur le graphe de dépendance, et employez-le pour démontrer que tout hypergraphe
k-uniforme dans lequel chaque arête en rencontre au plus $ 2^{k}/e - 1 $ autres admet une
2-coloration propre de ses sommets (aucune arête monochrome).
\item[(d)] Voici maintenant la piqûre. Chacune des démonstrations ci-dessus montre qu'une
coloration existe sans en produire aucune. Expliquez précisément quelle étape détruit la
construction, et lisez ce que l'on sait de sa récupération : Moser et Tardos (2010) ont
transformé le lemme local en algorithme, tandis que construire explicitement une coloration
de $ K_n $ sans clique monochrome de taille $ C\log n $ --- à la hauteur de ce qu'obtient
un tirage à pile ou face --- reste ouvert, les meilleures constructions explicites, après les
travaux de Chattopadhyay et Zuckerman de 2016 sur les extracteurs à deux sources, demeurant
très en deçà de la borne aléatoire.
\end{enumerate}

\end{problem}

\noindent La note d'Erdős de 1947 sur les nombres de Ramsey fait trois pages, et
l'argument de la question (a) en occupe un paragraphe ; elle a changé ce qu'un
combinatoricien entend par \og{} il existe \fg{}. Szele avait employé l'idée en 1943 pour produire
des tournois à nombreux chemins hamiltoniens, mais c'est Erdős qui en a fait une méthode, et
le livre d'Alon et Spencer qui en a fait un sujet. Philosophiquement, la méthode voisine avec
le tour à deux cas de PRF-45 et les contre-exemples invoqués par l'axiome du choix de
PRF-43 : une démonstration peut être étanche, quantitative, voire optimale, et refuser malgré
tout de livrer l'objet qu'elle promet. Ce refus s'est révélé fécond --- la recherche de
substituts constructifs a engendré la technique de compression d'entropie, la théorie des
extracteurs d'aléa et une bonne part de la dérandomisation.

\begin{refs}
\item Contexte : Méthode probabiliste --- Wikipédia (\url{https://fr.wikipedia.org/wiki/M%C3%A9thode_probabiliste})
\item Contexte : Lemme local de Lovász --- Wikipédia (\url{https://fr.wikipedia.org/wiki/Lemme_local_de_Lov%C3%A1sz})
\item Article : P. Erdős, \og{} Some remarks on the theory of graphs \fg{} (1947), \textit{Bulletin of the AMS} 53
\item Article : R. A. Moser \& G. Tardos, \og{} A constructive proof of the general Lovász Local Lemma \fg{} (2010), \textit{Journal of the ACM} 57
\item Lecture : N. Alon \& J. H. Spencer, \textit{The Probabilistic Method}
\end{refs}

\bigskip

\begin{problem}[{Une identité, deux démonstrations : bijection ou série génératrice --- Master, short}]
\textbf{PRF-49.} \textbf{Problème.} Démontrez le théorème d'Euler selon lequel le nombre
de partitions de n en parts impaires égale le nombre de partitions de n en parts distinctes
--- deux fois : une fois en exhibant une bijection explicite entre les deux ensembles de
partitions, une fois en démontrant l'identité de séries formelles
\[ \prod_{i \ge 1} \frac{1}{1 - q^{2i-1}} \;=\; \prod_{i \ge 1} \bigl(1 + q^{i}\bigr). \]
Décidez ensuite laquelle des deux vous emploieriez pour répondre à \og{} quelle partition de 100
en parts distinctes correspond à $ 99 + 1 $ ? \fg{}, et dites ce que cela vous apprend sur la
différence entre les deux.
\end{problem}

\noindent Euler a démontré l'identité analytiquement dans l'\textit{Introductio}
(1748) ; la bijection est venue plus tard, dans les travaux de Glaisher (1883), et procède en
scindant chaque part impaire répétée en puissances de deux --- une procédure, non une formule.
Le couple est le cas d'école standard d'un véritable clivage méthodologique : les séries
génératrices sont mécaniques, robustes, et se généralisent par l'algèbre, tandis qu'une
bijection est une pièce d'ingénierie combinatoire qui transporte réellement un objet vers un
autre et que l'on peut exécuter de tête. La démonstration par involution donnée par Franklin
en 1881 du théorème des nombres pentagonaux d'Euler est l'archétype du second genre, et la
même tension est à l'œuvre dans l'involution de PRF-47.

\begin{refs}
\item Contexte : Théorème de Glaisher --- Wikipédia (en anglais) (\url{https://en.wikipedia.org/wiki/Glaisher%27s_theorem})
\item Contexte : Partition d'un entier --- Wikipédia (\url{https://fr.wikipedia.org/wiki/Partition_d%27un_entier})
\item Lecture : G. E. Andrews, \textit{The Theory of Partitions} (1976), ch. 1
\end{refs}

\bigskip

\begin{problem}[{Le squelette de la démonstration de Bertrand par Erdős --- Master, short}]
\textbf{PRF-50.} \textbf{Problème.} Lisez la démonstration par Erdős du postulat de
Bertrand --- pour tout entier $ n \ge 1 $ il existe un nombre premier p avec
$ n <p \le 2n $ --- et reconstituez-en le squelette : dressez dans l'ordre la liste de
chaque lemme qu'elle emploie sur le coefficient binomial central $ \binom{2n}{n} $,
énoncez chacun comme une inégalité autonome, et signalez celui qui fait le vrai travail.
Déterminez ensuite quels lemmes il faudrait renforcer, et de combien, pour démontrer le même
énoncé sur l'intervalle plus court $ (n, \tfrac{3}{2}n] $.
\end{problem}

\noindent Bertrand a conjecturé l'énoncé en 1845 après l'avoir vérifié à la main
jusqu'à trois millions ; Tchebychev l'a démontré en 1852 avec une lourde machinerie ;
Ramanujan en a donné une démonstration plus courte en 1919 ; et en 1932 Erdős, âgé de
dix-neuf ans, a publié, comme tout premier article, un argument n'utilisant rien de plus que
la décomposition en facteurs premiers de $ \binom{2n}{n} $ et une majoration du produit des
nombres premiers. Démonter une démonstration achevée en ses pièces porteuses est un
savoir-faire distinct de celui d'en trouver une --- et le moyen de découvrir quelles hypothèses
étaient réellement nécessaires. Pièce jumelle des exercices de réécriture PRF-27 et
PRF-32.

\begin{refs}
\item Contexte : Postulat de Bertrand --- Wikipédia (\url{https://fr.wikipedia.org/wiki/Postulat_de_Bertrand})
\item Article : P. Erdős, \og{} Beweis eines Satzes von Tschebyschef \fg{} (1932), \textit{Acta Litterarum ac Scientiarum (Szeged)} 5, 194--198
\item Lecture : M. Aigner \& G. M. Ziegler, \textit{Proofs from THE BOOK}, le chapitre sur le postulat de Bertrand
\end{refs}

\bigskip

\begin{problem}[{Trouvez l'axiome du choix --- Master}]
\textbf{PRF-59.} \textbf{Problème.} Quelque part dans la plupart des arguments ci-dessous, une sélection
est faite une infinité de fois et personne ne le dit. L'exercice consiste à trouver l'étape,
à la nommer, et à déterminer ce qu'elle coûte réellement en axiome du choix.
\begin{enumerate}
\item[(a)] Pour chacune des démonstrations standard suivantes, localisez l'étape qui emploie une
forme de l'axiome du choix et identifiez le principe le plus faible qui suffise --- l'axiome
du choix complet, le choix dépendant, le choix dénombrable, ou aucun du tout : (i) tout
ensemble infini possède une partie infinie dénombrable ; (ii) une réunion dénombrable
d'ensembles dénombrables est dénombrable ; (iii) pour
$ f\colon \mathbb{R} \to \mathbb{R} $, la continuité en a au sens séquentiel entraîne la
continuité en a au sens $\varepsilon$--$\delta$ ; (iv) toute surjection admet un inverse à droite ; (v) tout
espace vectoriel possède une base.
\item[(b)] Deux des cinq ne sont pas de simples usagers d'un principe : ils lui sont équivalents.
Démontrez que \og{} toute surjection $ g\colon B \to A $ admet un inverse à droite \fg{} est
équivalent à l'axiome du choix au-dessus de ZF, et énoncez --- avec une référence, la
démonstration étant difficile --- le théorème de Blass selon lequel \og{} tout espace vectoriel
possède une base \fg{} lui est équivalent aussi.
\item[(c)] Exercez la discipline sur (iii) : réécrivez la démonstration habituelle de sorte que le
choix soit complètement explicite, en affichant la famille d'ensembles non vides et la
fonction de choix que l'on en extrait. Déterminez ensuite, à partir des références, si une
démonstration de (iii) sans axiome du choix est seulement possible, et dites quel genre
d'élément pourrait trancher une telle question.
\item[(d)] Enfin la méta-question, en un paragraphe. Gödel a démontré en 1938 que l'axiome du
choix est cohérent avec ZF, et Cohen a démontré en 1963 que sa négation l'est aussi : aucune
contradiction ne peut donc naître de son emploi. PRF-30 exigeait pourtant une démonstration
de Schröder--Bernstein sans axiome du choix. Pourquoi les mathématiciens se soucient-ils de
savoir quels axiomes une démonstration consomme ?
\end{enumerate}

\end{problem}

\noindent Zermelo a isolé l'axiome en 1904 afin de démontrer que tout ensemble peut
être bien ordonné, et la réaction fut féroce : Baire, Borel et Lebesgue ont publié des
objections en quelques mois, réunies dans les \textit{Cinq lettres sur la théorie des ensembles}
de Hadamard (1905), et la dispute portait exactement sur ce que la question (c) vous demande
de faire --- savoir si une démonstration peut invoquer une sélection qu'elle ne sait pas
décrire. L'image de Russell est celle dont tout le monde se souvient : choisir une chaussure
dans chacune d'une infinité de paires ne demande aucun axiome, car \og{} la gauche \fg{} est une
règle ; choisir une chaussette dans chacune d'une infinité de paires indiscernables demande
l'axiome. L'habitude de suivre l'axiome du choix à la trace n'est pas une manie d'antiquaire.
C'est la même habitude que les mathématiques à rebours ont transformée en programme de
recherche (PRF-63), et la même qui fait qu'un assistant de preuve vous demande quels
principes classiques vous voulez activer.

\begin{refs}
\item Contexte : Axiome du choix --- Wikipédia (\url{https://fr.wikipedia.org/wiki/Axiome_du_choix})
\item Contexte : Axiome du choix dénombrable --- Wikipédia (\url{https://fr.wikipedia.org/wiki/Axiome_du_choix_d%C3%A9nombrable})
\item Lecture : T. J. Jech, \textit{The Axiom of Choice} (North-Holland, 1973)
\item Article : A. Blass, \og{} Existence of bases implies the axiom of choice \fg{} (1984), \textit{Contemporary Mathematics} 31
\end{refs}

\bigskip

\begin{problem}[{Le Nullstellensatz combinatoire comme arme --- Master}]
\textbf{PRF-60.} \textbf{Problème.} Le Nullstellensatz combinatoire d'Alon. Soit F un corps et soit
$ f \in F[x_1, \dots, x_n] $ de degré total $ k_1 + \cdots + k_n $ ; supposons que le
coefficient de $ x_1^{k_1} \cdots x_n^{k_n} $ dans f soit non nul. Si
$ A_1, \dots, A_n $ sont des parties finies de F avec $ |A_i| > k_i $ pour tout i, alors
il existe des $ a_i \in A_i $ tels que $ f(a_1, \dots, a_n) \ne 0 $.
\begin{enumerate}
\item[(a)] Démontrez-le. (Route : le fait, à une variable, qu'un polynôme non nul de degré d a au
plus d racines, joint à la réduction de f modulo les polynômes
$ g_i(x_i) = \prod_{a \in A_i} (x_i - a) $, qui ne change aucun coefficient
pertinent.)
\item[(b)] Déduisez-en le théorème de Cauchy--Davenport : pour un nombre premier p et des parties
non vides A, B de $ \mathbb{Z}/p\mathbb{Z} $,
\[ |A + B| \;\ge\; \min\{\, p,\; |A| + |B| - 1 \,\}. \]
Tout l'art est de choisir le polynôme ; dites ensuite comment vous l'auriez
trouvé.
\item[(c)] Déduisez-en le théorème d'Erdős--Ginzburg--Ziv : parmi $ 2n - 1 $ entiers quelconques,
il en existe n dont la somme est divisible par n. Traitez d'abord le cas premier, puis
ramenez-y le cas général, et énoncez exactement où la réduction a besoin du cas premier et
non du seul énoncé.
\item[(d)] Vient ensuite la question méthodologique que cette page ne cesse de poser. Le
Nullstellensatz démontre qu'un bon choix existe en exhibant un coefficient non nul --- un
objet d'une tout autre nature que celui qui était promis. Comparez-le à la méthode
probabiliste de PRF-48 : laquelle des deux, si tant est que l'une des deux le puisse,
peut-on amener à produire un témoin pour (b), et à quel prix ?
\end{enumerate}

\end{problem}

\noindent La technique est née des travaux d'Alon et Tarsi de 1989 sur les
colorations et les orientations de graphes, a été affûtée avec Nathanson et Ruzsa en 1995-96,
et a reçu son nom et son énoncé net dans l'exposé de synthèse d'Alon de 1999, qui reste le
meilleur endroit où l'apprendre. Ce qui en a fait un jalon, c'est le transfert : des questions
additives sur des ensembles de résidus, des questions de coloration sur des graphes et des
questions de recouvrement sur le cube sont toutes devenues des questions portant sur un unique
coefficient d'un polynôme que l'on a le droit de concevoir. Ce geste --- encoder le problème
dans un polynôme, puis y lire un coefficient --- s'appelle aujourd'hui la méthode polynomiale,
et ses triomphes ultérieurs sont saisissants : la solution par Dvir, en 2008, du problème de
Kakeya sur les corps finis tient en une page, et l'idée de Croot, Lev et Pach a donné à
Ellenberg et Gijswijt, en 2016, la borne sur les cap sets, après des décennies de
stagnation.

\begin{refs}
\item Contexte : Nullstellensatz combinatoire --- Wikipédia (en anglais) (\url{https://en.wikipedia.org/wiki/Combinatorial_Nullstellensatz})
\item Contexte : Somme restreinte d'ensembles (théorème de Cauchy--Davenport) --- Wikipédia (\url{https://fr.wikipedia.org/wiki/Somme_restreinte_d%27ensembles})
\item Article : N. Alon, \og{} Combinatorial Nullstellensatz \fg{} (1999), \textit{Combinatorics, Probability and Computing} 8, 7--29
\item Lecture : L. Guth, \textit{Polynomial Methods in Combinatorics} (AMS, 2016)
\end{refs}

\bigskip

\begin{problem}[{Une seule diagonale, cinq théorèmes --- Master, short}]
\textbf{PRF-61.} \textbf{Problème.} Démontrez le cas ensembliste du théorème du point fixe
de Lawvere : si A et B sont des ensembles et si une application $ f\colon A \to B^{A} $ est
\textit{point-surjective} --- toute fonction $ A \to B $ est égale à $ f(a) $ pour un certain
$ a \in A $ --- alors toute $ g\colon B \to B $ possède un point fixe. Lisez ensuite la
contraposée comme une usine et faites-la tourner cinq fois, en en dérivant le théorème de
Cantor, le paradoxe de Russell, l'indécidabilité du problème de l'arrêt, le théorème
d'indéfinissabilité de Tarski et le premier théorème d'incomplétude de Gödel, et en disant à
chaque fois exactement ce que sont A, B, f et le g sans point fixe.
\end{problem}

\noindent PRF-21 vous demandait d'énoncer le canevas diagonal assez abstraitement
pour couvrir deux théorèmes ; voici ce qui arrive quand quelqu'un le fait pour de bon. William
Lawvere a publié le théorème en 1969, dans un article sur les catégories cartésiennes fermées
que la plupart des logiciens n'ont pas lu pendant trente ans, et il dit quelque chose de
véritablement surprenant : cinq résultats qui semblent appartenir à cinq disciplines
différentes sont un seul résultat plus cinq dictionnaires, et les \og{} paradoxes \fg{} ne sont que le
constat que certains g n'ont pas de point fixe. L'exposé de Yanofsky de 2003 en retire la
théorie des catégories et constitue la voie d'accès la plus rapide. Une fois cela fait,
remarquez ce qu'il en coûte : la démonstration unifiée explique pourquoi les théorèmes sont
vrais et n'est presque d'aucun secours pour démontrer quoi que ce soit de nouveau --- une
tension qu'il vaut la peine de garder à l'esprit.

\begin{refs}
\item Contexte : Théorème du point fixe de Lawvere --- Wikipédia (en anglais) (\url{https://en.wikipedia.org/wiki/Lawvere%27s_fixed-point_theorem})
\item Contexte : Lemme de diagonalisation --- Wikipédia (en anglais) (\url{https://en.wikipedia.org/wiki/Diagonal_lemma})
\item Article : F. W. Lawvere, \og{} Diagonal arguments and cartesian closed categories \fg{} (1969), \textit{Lecture Notes in Mathematics} 92
\item Article : N. S. Yanofsky, \og{} A universal approach to self-referential paradoxes, incompleteness and fixed points \fg{} (2003), \textit{Bulletin of Symbolic Logic} 9, 362--386, arXiv:math/0305282 (\url{https://arxiv.org/abs/math/0305282})
\end{refs}

\bigskip


\section*{Recherche}

\begin{problem}[{Après Bâle : les valeurs impaires de zêta --- Recherche}]
\textbf{PRF-33.} \textbf{Problème.} Euler ne s'est pas arrêté à $ \zeta(2) $ : il a démontré
que $ \zeta(2k) $ est un multiple rationnel de $ \pi^{2k} $ pour tout k.
\begin{enumerate}
\item[(a)] Échauffement :
étendez la comparaison de coefficients de PRF-16 pour démontrer
$ \zeta(4) = \sum 1/n^4 = \pi^4/90 $.
\item[(b)] Reconstituez, à partir du compte rendu informel de van der Poorten, la démonstration
donnée par Apéry en 1978 de l'irrationalité de $ \zeta(3) = \sum 1/n^3 $.
\item[(c)] Le problème ouvert : déterminer si $ \zeta(5) $ est irrationnel. \textbf{État :}
ouvert en 2026. Ce que l'on sait : $\zeta$(3) est irrationnel (Apéry, 1978) ; une infinité des
valeurs impaires $\zeta$(2k+1) sont irrationnelles (Ball--Rivoal, 2001) ; au moins l'une de $\zeta$(5),
$\zeta$(7), $\zeta$(9), $\zeta$(11) est irrationnelle (Zudilin, 2001). Aucune valeur impaire au-delà de $\zeta$(3)
n'a été tranchée, et la question de savoir si $\zeta$(3) est transcendant est ouverte elle
aussi.
\end{enumerate}

\end{problem}

\noindent Les descendants directs du problème de Bâle. Euler a cherché âprement une
forme close pour $\zeta$(3) et n'en a trouvé aucune ; nous croyons aujourd'hui qu'il n'en existe
pas au sens où il l'entendait. Quand Apéry, âgé de 61 ans, a annoncé en 1978 une
démonstration d'irrationalité n'utilisant que des mathématiques connues d'Euler, l'auditoire
s'est montré ouvertement incrédule --- le compte rendu de van der Poorten, \og{} A proof that
Euler missed\dots{} \fg{}, consigne la stupéfaction et reste la meilleure porte d'entrée. Le sujet mêle
aujourd'hui l'élémentaire et le transcendant : l'état de l'art ne parvient toujours pas à
trancher sur un nombre, $\zeta$(5), que n'importe quelle calculatrice imprime à mille
décimales.

\begin{refs}
\item Contexte : Théorème d'Apéry --- Wikipédia (\url{https://fr.wikipedia.org/wiki/Th%C3%A9or%C3%A8me_d%27Ap%C3%A9ry})
\item Article : A. van der Poorten, \og{} A proof that Euler missed\dots{} Apéry's proof of the irrationality of $\zeta$(3) \fg{} (1979), \textit{Mathematical Intelligencer} 1
\item Article : W. Zudilin, \og{} One of the numbers $\zeta$(5), $\zeta$(7), $\zeta$(9), $\zeta$(11) is irrational \fg{} (2001), \textit{Russian Mathematical Surveys} 56
\end{refs}

\bigskip

\begin{problem}[{Les nombres premiers de Wilson --- Recherche}]
\textbf{PRF-34.} \textbf{Problème.} Le théorème de Wilson (PRF-20) dit que p divise
$ (p-1)! + 1 $ pour tout nombre premier p. Un \textit{nombre premier de Wilson} est un
nombre premier pour lequel la divisibilité est doublement forte :
$ p^2 \mid (p-1)! + 1 $.
\begin{enumerate}
\item[(a)] Vérifiez à la main que 5
et 13 sont des nombres premiers de Wilson, et à l'ordinateur que 563 en est un.
\item[(b)] Le problème ouvert : démontrez
ou réfutez qu'il existe une infinité de nombres premiers de Wilson --- ou même qu'il en existe
un quatrième. \textbf{État :} ouvert en 2026. On n'en connaît que 5, 13 et 563 ; une recherche
exhaustive de Costa, Gerbicz et Harvey n'en a trouvé aucun autre en dessous de
$ 2 \times 10^{13} $.
\item[(c)] Expliquez l'heuristique standard --- traiter $ ((p-1)!+1)/p \bmod p $ comme un résidu
aléatoire --- et déduisez-en la prédiction selon laquelle le nombre de nombres premiers de
Wilson jusqu'à x devrait croître comme log log x. Que faudrait-il pour transformer une telle
heuristique en démonstration ?
\end{enumerate}

\end{problem}

\noindent Un spécimen parfait de frontière de recherche poussant directement sur un
théorème classique : renforcez une congruence d'une puissance de p et la démontrabilité
s'effondre. L'heuristique de la question (c) est crue de tous et démontrée par personne --- le
même mur que celui qui se dresse devant les nombres premiers de Fermat, les nombres premiers
de Mersenne et la conjecture des nombres premiers jumeaux (voyez la page des problèmes
ouverts). La recherche de Costa, Gerbicz et Harvey est aussi un jalon de la théorie
algorithmique des nombres : même \textit{vérifier} la congruence de Wilson assez vite pour
atteindre $ 2 \times 10^{13} $ a exigé des idées nouvelles sur le calcul des factorielles
modulo p.

\begin{refs}
\item Contexte : Nombre premier de Wilson --- Wikipédia (\url{https://fr.wikipedia.org/wiki/Nombre_premier_de_Wilson})
\item Article : E. Costa, R. Gerbicz \& D. Harvey, \og{} A search for Wilson primes \fg{} (2014), \textit{Mathematics of Computation} 83, arXiv:1209.3436 (\url{https://arxiv.org/abs/1209.3436})
\end{refs}

\bigskip

\begin{problem}[{Toutes les tautologies ont-elles des démonstrations courtes ? --- Recherche}]
\textbf{PRF-35.} \textbf{Problème.} Un système de démonstration propositionnel est, informellement, tout
format correct et complet permettant de certifier qu'une formule booléenne est une
tautologie, avec des certificats vérifiables en temps polynomial.
\begin{enumerate}
\item[(a)] À partir des références, travaillez le
théorème de Cook--Reckhow : NP = coNP si et seulement s'il existe un système de démonstration
propositionnel admettant des démonstrations de taille polynomiale pour toutes les
tautologies. Vérifiez vous-même le sens facile.
\item[(b)] Étudiez le théorème de Haken de 1985 : le système de démonstration par résolution exige
des démonstrations de taille exponentielle pour les tautologies des tiroirs --- l'encodage
propositionnel du principe de PRF-07 selon lequel n+1 objets ne peuvent entrer
injectivement dans n tiroirs.
\item[(c)] Le problème ouvert : exhiber une
famille de tautologies exigeant des démonstrations de taille superpolynomiale dans un
système de Frege (la logique propositionnelle ordinaire des manuels) --- ou démontrer qu'il
n'en existe aucune. \textbf{État :} ouvert en 2026 ; on ne connaît aucune minoration
superpolynomiale pour Frege, et la question NP contre coNP reste ouverte (une réponse
négative impliquerait P $\ne$ NP).
\end{enumerate}

\end{problem}

\noindent Ici, la démonstration elle-même devient l'objet des mathématiques. Cook et
Reckhow ont fondé la complexité des preuves en 1979 précisément pour attaquer P contre NP par
le côté logique, et la minoration de Haken --- le fait \og{} évident \fg{} des tiroirs n'admet pas de
démonstration courte par résolution --- fut la première grande victoire du programme, avec une
morale sur laquelle tout rédacteur de démonstrations devrait méditer : un énoncé peut être
facile à croire, facile à démontrer avec la bonne idée (une ligne de dénombrement !), et
démontrablement impossible à démontrer brièvement dans un langage restreint. La distance qui
sépare la résolution des systèmes de Frege a résisté à quarante ans d'efforts. Voyez la page
des problèmes ouverts pour l'histoire plus large de P contre NP.

\begin{refs}
\item Contexte : Complexité des preuves --- Wikipédia (\url{https://fr.wikipedia.org/wiki/Complexit%C3%A9_des_preuves})
\item Article : S. A. Cook \& R. A. Reckhow, \og{} The relative efficiency of propositional proof systems \fg{} (1979), \textit{Journal of Symbolic Logic} 44
\item Article : A. Haken, \og{} The intractability of resolution \fg{} (1985), \textit{Theoretical Computer Science} 39
\end{refs}

\bigskip

\begin{problem}[{Les droites ordinaires après Sylvester--Gallai --- Recherche}]
\textbf{PRF-36.} \textbf{Problème.} Sylvester--Gallai (PRF-26) garantit une droite ordinaire ; combien
doit-il y en avoir ? La conjecture de Dirac--Motzkin affirme que n points non alignés
déterminent au moins $ n/2 $ droites ordinaires lorsque n est pair (avec des
configurations exceptionnelles connues pour certains petits n impairs).
\begin{enumerate}
\item[(a)] Construisez, pour une infinité de n pairs, une
configuration atteignant exactement $ n/2 $ droites ordinaires.
\item[(b)] Trouvez une configuration de
7 points n'ayant que 3 droites ordinaires.
\item[(c)] Étudiez le théorème de Green et Tao (2013) : pour tout
n assez grand, tout ensemble de n points non alignés détermine au moins $ n/2 $
droites ordinaires. \textbf{État :} résolu pour tout n assez grand par Green et Tao
(2013), dont le seuil \og{} assez grand \fg{} est astronomiquement grand et ineffectif en
pratique ; obtenir un seuil explicite et à taille humaine --- et donc la conjecture pour
tout n --- reste ouvert en 2026.
\end{enumerate}

\end{problem}

\noindent Un arc long d'un siècle : le casse-tête de Sylvester en 1893, la
démonstration de Gallai en 1944, les deux lignes parfaites de Kelly, puis un saut de
difficulté d'un tout autre ordre --- la démonstration de Green et Tao fait intervenir la
géométrie algébrique des courbes cubiques, montrant que les configurations à peu de droites
ordinaires doivent épouser une cubique (la configuration hessienne de PRF-26 n'est pas un
hasard). C'est un modèle de la façon dont un théorème classique résolu continue d'engendrer
de la recherche : remplacez \og{} au moins une \fg{} par \og{} combien \fg{} et un sujet élémentaire rencontre
la frontière. Leur article a du même coup réglé, pour n grand, le problème de la plantation
du verger de 1868.

\begin{refs}
\item Contexte : Théorème de Sylvester-Gallai --- Wikipédia (\url{https://fr.wikipedia.org/wiki/Th%C3%A9or%C3%A8me_de_Sylvester-Gallai})
\item Article : B. Green \& T. Tao, \og{} On sets defining few ordinary lines \fg{} (2013), \textit{Discrete \& Computational Geometry} 50, arXiv:1208.4714 (\url{https://arxiv.org/abs/1208.4714})
\end{refs}

\bigskip

\begin{problem}[{Quand le rapporteur est une machine --- Recherche}]
\textbf{PRF-51.} \textbf{Problème.} Plusieurs théorèmes reposent désormais sur des démonstrations
qu'aucun humain n'a lues intégralement : le théorème des quatre couleurs (Appel--Haken,
1976 ; simplifié par Robertson, Sanders, Seymour et Thomas, 1997 ; formalisé en Coq par
Gonthier, 2005), la conjecture de Kepler (Hales, 1998, dont les rapporteurs ont refusé de la
certifier complètement --- formalisée par le projet Flyspeck, achevé en 2014 et publié en
2017), et le théorème de Feit--Thompson sur l'ordre impair (formalisé en Coq par Gonthier et
ses collaborateurs, 2012).
\begin{enumerate}
\item[(a)] Énoncez exactement ce qu'établit un assistant de preuve. Décrivez la base de confiance
d'une démonstration en Lean ou en Coq --- noyau, axiomes, l'énoncé du théorème lui-même --- et
expliquez pourquoi \og{} un ordinateur l'a vérifié \fg{} est une affirmation portant sur un système
formel donné, et pourquoi l'énoncé que l'on formalise est justement la partie qu'une machine
ne peut pas vérifier à votre place.
\item[(b)] Formalisez, en Lean 4 avec mathlib, un théorème que vous savez démontrer en une page :
l'irrationalité de $ \sqrt{2} $ (PRF-02) ou l'infinité des nombres premiers (PRF-14).
Mesurez le rapport de la longueur formelle à la longueur informelle --- le \og{} facteur de de
Bruijn \fg{} --- et identifiez quelles étapes informelles se sont le plus dilatées. Débattez de ce
que révèle cette dilatation.
\item[(c)] Prenez position, éléments à l'appui, sur la question de la vérifiabilité humaine
soulevée par Tymoczko en 1979 : si personne ne peut vérifier une démonstration, en quel sens
la communauté connaît-elle le théorème ? Opposez le cas du théorème des quatre couleurs à
celui du Liquid Tensor Experiment (2020-2022), où Scholze a demandé la vérification par
machine d'une démonstration qu'il avait écrite et à laquelle il ne se fiait pas
entièrement, et à celui du théorème polynomial de Freiman--Ruzsa de Gowers, Green, Manners et
Tao (2023), formalisé quelques semaines après son annonce.
\textbf{État :} ouvert, et en partie non mathématique. On ne connaît aucune démonstration sans
ordinateur du théorème des quatre couleurs ni de la conjecture de Kepler, et aucun théorème
ne dit qu'il ne peut en exister ; savoir si les démonstrations vérifiées par machine
resteront une spécialité ou deviendront la norme se décide dans la pratique, à ce jour en
2026.
\end{enumerate}

\end{problem}

\noindent La démonstration du théorème des quatre couleurs par Appel et Haken en
1976 a vérifié 1 834 configurations réductibles par ordinateur et a provoqué une véritable
crise : des mathématiciens qui avaient passé un siècle à échouer à démontrer le théorème
s'entendaient dire qu'il était vrai par une machine dont ils ne pouvaient auditer la sortie.
Le développement en Coq de Gonthier, en 2005, a répondu à la forme la plus étroite de
l'objection --- l'argument est réellement correct, jusqu'aux axiomes --- sans toucher à la plus
large, à savoir que la compréhension n'a jamais été obtenue. Trente ans plus tard, l'humeur
s'est inversée : Hales s'est tourné vers la formalisation parce que des rapporteurs humains
ne parvenaient pas à finir, et Scholze parce qu'il voulait faire vérifier une démonstration
de son cru. La question de cette section n'est donc pas \og{} les ordinateurs auront-ils raison \fg{}
mais à quoi servent les mathématiques ; voyez aussi PRF-32 sur l'écriture d'un même théorème
pour trois publics, et PRF-35 sur les démonstrations qu'un système formel ne peut exprimer
qu'à une longueur déraisonnable.

\begin{refs}
\item Contexte : Preuve assistée par ordinateur --- Wikipédia (\url{https://fr.wikipedia.org/wiki/Preuve_assist%C3%A9e_par_ordinateur})
\item Contexte : Théorème des quatre couleurs --- Wikipédia (\url{https://fr.wikipedia.org/wiki/Th%C3%A9or%C3%A8me_des_quatre_couleurs})
\item Contexte : Lean (assistant de preuve) --- Wikipédia (\url{https://fr.wikipedia.org/wiki/Lean_%28assistant_de_preuve%29})
\item Article : G. Gonthier, \og{} Formal Proof --- The Four-Color Theorem \fg{} (2008), \textit{Notices of the AMS} 55
\item Article : T. C. Hales et al., \og{} A formal proof of the Kepler conjecture \fg{} (2017), \textit{Forum of Mathematics, Pi} 5
\end{refs}

\bigskip

\begin{problem}[{À quelle vitesse croissent les nombres de Ramsey diagonaux ? --- Recherche, short}]
\textbf{PRF-52.} \textbf{Problème.} Erdős--Szekeres (1935) donne
$ R(k,k) \le 4^{k} $ et la coloration aléatoire de PRF-48 donne
$ R(k,k) >2^{k/2} $ ; en 2023 Campos, Griffiths, Morris et Sahasrabudhe ont démontré
$ R(k,k) \le (4-\varepsilon)^{k} $ pour une constante absolue $ \varepsilon >0 $,
la première amélioration exponentielle en quatre-vingt-huit ans. Déterminez le véritable taux
de croissance exponentielle : décidez si $ \lim_{k\to\infty} R(k,k)^{1/k} $ existe et, si
oui, calculez-le.
\end{problem}

\noindent \textbf{État : ouvert} --- même l'existence de la limite est inconnue, et
Erdős avait notoirement offert des prix distincts pour les deux moitiés de la question. Lisez
la démonstration de 2023 pour ce qu'elle fait à la constante, et pas seulement pour la
constante : l'argument est un tour de force de comptabilité bâti sur une construction
algorithmique dite \og{} du livre \fg{}, et il montre que huit décennies d'immobilité tenaient à une
idée manquante plutôt qu'à une barrière. La minoration a elle aussi bougé récemment, la
constante de l'estimation d'Erdős de 1947 ayant été améliorée pour la première fois.

\begin{refs}
\item Contexte : Théorème de Ramsey --- Wikipédia (\url{https://fr.wikipedia.org/wiki/Th%C3%A9or%C3%A8me_de_Ramsey})
\item Article : M. Campos, S. Griffiths, R. Morris \& J. Sahasrabudhe, \og{} An exponential improvement for diagonal Ramsey \fg{} (2023), \textit{Annals of Mathematics} 203 (2026), arXiv:2303.09521 (\url{https://arxiv.org/abs/2303.09521})
\end{refs}

\bigskip

\begin{problem}[{La conjecture de Frankl sur les familles stables par union --- Recherche, short}]
\textbf{PRF-53.} \textbf{Problème.} La conjecture de Frankl (1979) affirme que dans toute
famille finie d'ensembles stable par union --- autre que la famille $ \{\emptyset\} $ --- un
élément appartient à au moins la moitié des ensembles. Lisez l'argument de Gilmer de 2022,
qui emploie l'entropie de Shannon pour démontrer qu'un élément appartient à au moins une
fraction fixe et strictement positive des ensembles, isolez l'étape exacte où sa constante
est produite, et déterminez si la méthode peut être poussée jusqu'à $ 1/2 $.
\end{problem}

\noindent \textbf{État : ouvert.} Pendant quarante-trois ans on n'a rien su au-delà de
cas particuliers --- pas même qu'\textit{une} fraction constante convienne --- puis un court
argument a importé une technique de théorie de l'information venue d'un tout autre quartier
de la combinatoire. La constante de Gilmer, d'environ $ 0.01 $, a été améliorée en quelques
semaines à $ (3-\sqrt5)/2 \approx 0.38197 $ par plusieurs groupes indépendants, et des
raffinements ultérieurs l'ont poussée un peu au-delà de $ 0.382 $ ; savoir si la voie
entropique peut seulement atteindre $ 1/2 $ est en soi non tranché. Une démonstration en
direct que c'est une technique de démonstration, et non simplement une démonstration, qui est
la ressource rare.

\begin{refs}
\item Contexte : Conjecture des familles stables par unions --- Wikipédia (\url{https://fr.wikipedia.org/wiki/Conjecture_des_familles_stables_par_unions})
\item Article : J. Gilmer, \og{} A constant lower bound for the union-closed sets conjecture \fg{} (2022), arXiv:2211.09055 (\url{https://arxiv.org/abs/2211.09055})
\end{refs}

\bigskip

\begin{problem}[{Comment lire un article --- Recherche}]
\textbf{PRF-62.} \textbf{Problème.} Lire la recherche est un savoir-faire que personne n'enseigne et que
tout le monde est censé posséder. Choisissez un article dans cette liste --- tous courts, tous
réels, tous cités ailleurs sur cette page : Erdős, \og{} Some remarks on the theory of graphs \fg{}
(1947) ; Furstenberg, \og{} On the infinitude of primes \fg{} (1955) ; la démonstration d'une seule
phrase de Zagier (1990) ; Gilmer, \og{} A constant lower bound for the union-closed sets
conjecture \fg{} (2022). Effectuez ensuite les passes ci-dessous dans l'ordre, en écrivant votre
production à chaque étape avant de passer à la suivante.
\begin{enumerate}
\item[(a)] Première passe, cinq minutes : titre, résumé, introduction, titres de sections,
conclusions, références. Sans lire plus avant, écrivez ce que vous croyez que l'article
démontre, ce qu'il suppose, avec les travaux de qui il est en concurrence, et quelles trois
références il vous faudrait connaître pour le suivre.
\item[(b)] Deuxième passe, une heure : lisez le corps du texte, mais sautez les démonstrations.
Produisez un graphe de dépendances des énoncés numérotés de l'article --- quel lemme alimente
quel théorème --- et signalez chaque résultat importé de l'extérieur de l'article.
\item[(c)] Troisième passe : choisissez le lemme que votre graphe désigne comme le plus porteur, et
essayez de le démontrer vous-même à partir de son seul énoncé, avant de lire la
démonstration de l'auteur. Notez jusqu'où vous êtes allé et où vous vous êtes arrêté ; cette
frontière est l'apport réel de l'article pour vous.
\item[(d)] Rédigez un rapport de rapporteur : le théorème principal dans vos propres mots, l'idée
nouvelle, l'étape que vous n'avez pas su vérifier, trois questions à l'auteur, et une
recommandation. Comparez ensuite votre rapport à la réception de l'article dans la
littérature.
\item[(e)] Enfin, la passe historique. Trouvez l'article auquel celui-ci répond et le premier
article qui le cite, et dites en un paragraphe ce qui a changé entre les deux.
\end{enumerate}

\end{problem}

\noindent Un article de recherche n'est pas écrit dans l'ordre où il a été
découvert ; l'introduction est composée en dernier, les tentatives ratées sont supprimées, et
le tout est comprimé pour des lecteurs qui partagent déjà un énorme corps de contexte. Les
premières lectures donnent donc à tout le monde, et durablement, un sentiment d'échec, et le
remède est un protocole plutôt qu'un surcroît d'effort --- la méthode des trois passes de
Keshav, écrite pour l'informatique en 2007, est la plus connue et se transpose sans
modification. L'essai de Thurston explique pourquoi cette compression a lieu et ce qu'elle
fait perdre : les mathématiques circulent entre les personnes sous forme de compréhension,
dont l'article est un encodage avec pertes. La question (d) n'est pas une simulation.
L'évaluation par les pairs consiste en des gens qui font exactement cela, le plus souvent
anonymement et le plus souvent gratuitement, et le moyen le plus rapide de comprendre comment
la vérité mathématique est réellement certifiée est d'en rédiger un honnêtement. Comparez avec
PRF-50, qui démonte une seule démonstration achevée, et avec PRF-32, qui en rebâtit une pour
trois publics.

\begin{refs}
\item Contexte : Évaluation par les pairs --- Wikipédia (\url{https://fr.wikipedia.org/wiki/%C3%89valuation_par_les_pairs})
\item Contexte : arXiv --- Wikipédia (\url{https://fr.wikipedia.org/wiki/ArXiv})
\item Article : S. Keshav, \og{} How to Read a Paper \fg{} (2007), \textit{ACM SIGCOMM Computer Communication Review} 37(3), PDF (en anglais) (\url{https://web.stanford.edu/class/ee384m/Handouts/HowtoReadPaper.pdf})
\item Article : W. P. Thurston, \og{} On Proof and Progress in Mathematics \fg{} (1994), arXiv:math/9404236 (\url{https://arxiv.org/abs/math/9404236})
\end{refs}

\bigskip

\begin{problem}[{Combien peu coûte un théorème ? --- Recherche, short}]
\textbf{PRF-63.} \textbf{Problème.} L'arithmétique des fonctions élémentaires (EFA) est le
fragment de l'arithmétique doté de l'addition, de la multiplication et de l'exponentiation,
dont la récurrence est restreinte aux formules à quantificateurs bornés --- bien plus faible
que l'arithmétique de Peano. La grande conjecture de Harvey Friedman (1999) affirme que tout
théorème publié dans les \textit{Annals of Mathematics} dont l'énoncé ne mentionne que des
objets finitaires peut être démontré dans EFA : tranchez-la, ou réglez le seul cas d'école du
grand théorème de Fermat.
\end{problem}

\noindent \textbf{État : ouvert}, et pas du genre qu'une quinzaine astucieuse
refermera, puisqu'il s'agit d'une affirmation portant sur toute une littérature publiée : une
réfutation suppose d'exhiber un théorème arithmétique des \textit{Annals} dont on démontre
qu'il exige plus qu'EFA, et une démonstration suppose une réduction générale que personne ne
sait actuellement imaginer. Friedman a posté la conjecture sur la liste de diffusion
Foundations of Mathematics le 16 avril 1999. Le cas d'école de Fermat est concret et
instructif : la démonstration de Wiles, telle qu'elle est écrite, passe par la cohomologie
étale de Grothendieck, que l'on met normalement en place avec des univers --- hypothèse bien
au-delà de ZFC --- et Colin McLarty a montré que la machinerie réellement nécessaire peut être
rebâtie dans l'arithmétique d'ordre fini, laquelle est énormément plus faible que les univers
et encore énormément plus forte qu'EFA. Derrière la conjecture se tient la découverte
empirique centrale des mathématiques à rebours : les théorèmes ordinaires se regroupent en une
poignée de niveaux de force, de sorte que les mathématiciens en exercice semblent employer
bien moins de puissance axiomatique qu'ils n'y auraient droit, et personne ne sait
pourquoi.

\begin{refs}
\item Contexte : Arithmétique des fonctions élémentaires --- Wikipédia (en anglais) (\url{https://en.wikipedia.org/wiki/Elementary_function_arithmetic})
\item Contexte : Mathématiques à rebours --- Wikipédia (\url{https://fr.wikipedia.org/wiki/Math%C3%A9matiques_%C3%A0_rebours})
\item Article : C. McLarty, \og{} What does it take to prove Fermat's Last Theorem? Grothendieck and the logic of number theory \fg{} (2010), \textit{Bulletin of Symbolic Logic} 16
\item Lecture : S. G. Simpson, \textit{Subsystems of Second Order Arithmetic} (2e éd., Cambridge, 2009)
\end{refs}

\bigskip


\section*{Pour aller plus loin}


\vfill
\noindent\emph{Hors de la Caverne} --- des probl\`emes, pas de r\'eponses.
\end{document}
